\documentclass[12pt]{amsart}

\usepackage{amscd,amstext,amsfonts,amsthm,amsxtra,amsmath}
\usepackage{amssymb}

\usepackage{mathrsfs}
 \usepackage[all]{xy}

\usepackage[margin=1in]{geometry}
\oddsidemargin=0pt \evensidemargin=0pt
\topmargin=0in
\setlength{\textwidth}{6.3 in}
%\usepackage[top=29truemm,bottom=23truemm,left=27truemm,right=27truemm]{geometry}


\theoremstyle{plain}
\newtheorem{theorem} {Theorem} [section]
\newtheorem{lemma} [theorem]{Lemma}
\newtheorem{proposition}[theorem]{Proposition}
\newtheorem{corollary} [theorem]{Corollary}
%
\theoremstyle{definition}
\newtheorem{definition} [theorem] {Definition}
\newtheorem{example} [theorem] {Example}
\newtheorem{notation} [theorem]  {Notation}
\newtheorem{remark}[theorem]{Remark}

\pagenumbering{arabic}

\def\Sym{\operatorname{Sym}}
\def\Gr{\operatorname{Gr}}
\def\GL{\operatorname{GL}}
\def\aut{\mathfrak{aut}}
\def\Hom{\operatorname{Hom}}
\def\Ker{\operatorname{Ker}}
\def\Symb{\operatorname{Symb}}
\def\Isom{\operatorname{Isom}}
\def\Aut{\operatorname{Aut}}
\def\Sing{\operatorname{Sing}}
\def\span{\operatorname{span}}

\title[Characterizations of horospherical varieties]{Characterizations of smooth projective horospherical varieties of Picard number one}

\date{\today}

\author[J. Hong]{Jaehyun Hong}
\address{Center for Complex Geometry,  Institute for Basic Science (IBS), 55  Expo-ro, Yuseong-gu, Daejeon, Korea 34126}
\email{jhhong00@ibs.re.kr}

\author[S.-Y. Kim]{Shin-young Kim}
  \address{Center for Geometry and Physics, Institute for Basic Science (IBS), 77 Cheongam-ro, Nam-gu, Pohang, Gyeongbuk, Korea 37673}
\email{shinyoungkim@ibs.re.kr}

\keywords{geometric structure, local equivalence, horospherical variety, Cartan geometry}
\subjclass[2000]{32M12, 53A55, 14J45}

%\linenumbers

\begin{document}

\begin{abstract}
Let $X$ be a  smooth projective horospherical variety of Picard number one, of type $B_m$ for $m \geq 3$ or of type $F_4$. We show that any uniruled projective manifold with Picard number one is biholomorphic to $X$ if the variety of minimal rational tangents at a general point is projectively equivalent to that of $X$. %By a Cartan connection from the methods of W-normal complete step prolongation, it enough to show the vanishing of Lie algebra cohomologies up to second ones to get a local flatness of the geometric structure arising from the variety of minimal rational tangents.
\end{abstract}



\maketitle
%\tableofcontents

\section{Introductions}



The theory of varieties of minimal rational tangents, introduced by Hwang and Mok, has played an important role in the complex geometry of a uniruled projective manifold of Picard number one, such as the deformation rigidity, the stability of the tangent bundles,
and the extension property. The general philosophy in this theory is that one should be able to recover the complex geometry of a uniruled projective manifold of Picard number
one from the projective geometry of its varieties of minimal rational tangents. There are many results manifesting this philosophy, one of which is Cartan-Fubini extension Theorem (Theorem 1.2 of \cite{HM01}). This allows us to extend biholomorphisms between connected (Euclidean) open subsets of two Fano manifolds of Picard number one which preserve varieties of minimal rational tangents under mild geometric conditions. Another result in this direction is recognizing rational homogeneous varieties of Picard number one by their varieties of minimal rational tangents.

\begin{theorem} [Main Theorem of \cite{Mk08}, Main Theorem of \cite{HH}] \label{previous results} Let $X$ be a rational homogeneous variety associated to a long root and let $\mathcal C_o(X) \subset \mathbb PT_{o}X$ be the variety of minimal rational tangents at a base point $o \in X$. Let $M$ be a Fano manifold of Picard number one and $\mathcal C_x(M) \subset \mathbb PT_xM$ be the variety of minimal rational tangents at a general point $x \in X$ associated to a minimal dominant rational component. Suppose that $\mathcal C_x(M) \subset \mathbb PT_xM$ are projectively equivalent to $\mathcal C_o (X)\subset \mathbb PT_oX$ for general $x \in M$. Then $M$ is biholomorphic to $X$.
\end{theorem}

Theorem \ref{previous results} follows from the existence of a biholomorphism $\phi$ from a connected open subset $U$ of $M$ onto an open subset of $X$ whose differential $d\phi$ maps $\mathcal C_x(M)$ onto $\mathcal C_{\phi(x)}(X)$ for any $x \in U$. One  of the effective ways to solve this kind of local equivalence problems is  constructing  Cartan connections. In the  cases dealt with in Theorem \ref{previous results}, the existence of Cartan connections solving local equivalence problems associated to the subbundle $\mathcal C(X) \subset \mathbb P(TX)$  was proved by  Tanaka (\cite{Ta}).

More generally, given a quasi-homogeneous variety $X$ which is uniruled and smooth,  we may ask the same question, whether we can characterize $X$, or, the open dense orbit $X^0$ of the automorphism group of $X$, in terms of the variety of minimal rational tangents.    In order to characterize the open dense orbit $X^0$ of the automorphism group of $X$ in term of the variety of minimal rational tangents, it is natural to assume that there is a minimal rational curve in $X$ contained in $X^0$ completely. This happens, for example, when the boundary $X \backslash X^0$ has codimension $\geq 2$ in $X$. Among such quasi-homogeneous varieties, we will consider smooth projective horospherical varieties of Picard number one.



For a reductive   group $L$, a normal $L$-variety $X$ is said to be \emph{horospherical} if it has an open $L$-orbit $L/H$ whose isotropy group $H$ contains the unipotent part of a Borel subgroup of $L$. Then the normalizer $P$ of $H$ in $L$ is a parabolic subgroup of $L$ and the open orbit $L/H$ is isomorphic to a $(\mathbb C^*)^r$- bundle over the rational homogeneous variety $L/P$ for some $r \in \mathbb Z_{\geq 0}$. %In this case, $r$ is called the {\it rank} of $X$.

  Pasquier classified smooth horospherical varieties of Picard number one and obtained the following result.

\begin{theorem}[Theorem 0.1  of \cite{Pa}]%\label{Horospheical}
Let $L$ be a reductive group.
 Let $X$ be a smooth nonhomogeneous projective horospherical $L$-variety with Picard number one.
 Then $X$ is uniquely determined by its two closed $L$-orbits $ Y$ and $Z$, isomorphic to $L/P_Y$ and $L/P_Z$, respectively; and $(L, P_Y, P_Z)$ in one of the triples of the following list.
 \begin{enumerate}
   \item[\rm(1)] $(B_m,{\alpha_{m-1}},{\alpha_{m}})$ for $m\geq3$;
   \item[\rm(2)] $(B_3,{\alpha_{1}},{\alpha_{3}})$;
   \item[\rm(3)] $(C_m,{\alpha_{i+1}},{\alpha_{i}})$ for $m\geq2$, $i\in \{1,\ldots, m-1\}$;
   \item[\rm(4)] $(F_4,{\alpha_{2}},{\alpha_3})$;
   \item[\rm(5)] $(G_2,{\alpha_{2}},{\alpha_1})$.


 \end{enumerate}

% Furthermore, $\Aut(X)$ has two orbits in $X$,   $Z$ and the complement of $Z$ in $X$.
\end{theorem}

%......
The main result in this paper is to characterize smooth horospherical varieties of type $(B_m, \alpha_{m-1}, \alpha_m)$ or of type $(F_4, \alpha_2, \alpha_3)$ by using the variety of minimal rational tangents as Theorem \ref{previous results}.

\begin{theorem} \label{main results}
Let $X$ be a smooth projective horospherical variety of Picard number one. Assume that $X$ is of type $(B_m, \alpha_{m-1}, \alpha_m)$ for $m \geq 3$ or of type $(F_4, \alpha_2, \alpha_3)$.
Let $\mathcal C_o(X)$ denote the variety of minimal rational tangents of $X$ at a base point $o$ of $X$.
Let $M$ be a Fano manifold of Picard number one and $\mathcal C_x(M) \subset \mathbb PT_xM$ be the variety of minimal rational tangents at a general point $x \in X$ associated to a minimal dominant rational component. Suppose that $\mathcal C_x(M) \subset \mathbb PT_xM$ are projectively equivalent to $\mathcal C_o (X)\subset \mathbb PT_oX$ for general $x \in M$. Then $M$ is biholomorphic to $X$.
\end{theorem}

 The characterization problem for the case when $X$ is of type $(G_2, \alpha_2, \alpha_1)$ was discussed in \cite{Hw15}. Most of the arguments in \cite{Hw15}    can also be applied to the case when $X$ is of type $(B_m, \alpha_{m-1}, \alpha_m)$ or of type $(F_4, \alpha_2, \alpha_3)$, which is done in Section \ref{sect:horospherical BF type}.
%
 For the   geometric structures associated to them,
 the second author suggested a way to construct a Cartan connection  which solves   the local equivalence problem by showing that the condition (C) in \cite{Mo93} is satisfied
(Proposition 53 of \cite{Ki}).
Due to a gap in the proof of Theorem 17 of \cite{Ki}, instead,  we use the prolongation methods developed in \cite{HM}.

For the smooth horospherical variety of type $(C_m, \alpha_{i+1}, \alpha_i)$ or of type $(G_2, \alpha_2, \alpha_1)$,  Hwang and   Li  solve the characterization problem,  proving   that these horospherical varieties can be characterized by their varieties of minimal rational tangents (\cite{HL19}, \cite{HL21}).  They use  Theorem 2.16 of \cite{HL19}  to solve the local equivalence problem.
A new input of this paper is the use of the prolongation methods   and the computation of the Lie algebra cohomology, which replace  the role of Theorem 2.16 of \cite{HL19}.
As a result, we don't need to show the vanishing of holomorphic sections of all types but  the vanishing of only those associated with the nonvanishing Lie algebra cohomology. 

The computation of Lie algebra cohomology $H^2(\mathfrak g_-, \Gamma)$ is of independent interest, where $\mathfrak g_-$ is the negative part of a graded Lie algebra $\frak g$ and $\Gamma$ is a representation of $\frak g$. When $\frak g$ is not semisimple, we cannot apply the theory of Kostant directly to compute it (\cite{Ko61}). We suggest a way to compute the Lie algebra cohomology $H^2(\mathfrak g_-, \Gamma)$ by reducing it to the computation of  the Lie algebra cohomology associated with the restriction of the action to the semisimple part of $\frak g$. It is worth comparing our procedure with the Lyndon-Hochschild-Serre spectral sequence, which is a tool to compute the Lie algebra cohomology when $\frak g$ is not semisimple (\cite{HS53}, \cite{HS73}). \\


This paper is organized as follows.
We review the theory of varieties of minimal rational tangents and collect results on the second fundamental forms and the third fundamental forms of varieties of minimal rational tangents in Section \ref{sect:vmrt}. Section \ref{sect:geometric structures} devotes to the theory of geometric structures, correspondence between $G_0$-structures and ${\bf S}$-structures, and we review the prolongation methods which gives  the local  equivalence under the vanishing of sections of vector bundles associated to Lie algebra cohomology of degree 2. We describe the varieties of minimal rational tangents of smooth horospherical varieties of Picard number one in Section  \ref{horospherical varieties}.
In Section \ref{sect:horospherical BF type} -- Section \ref{sect:H2 cohomology},   restricting ourselves to horospherical varieties of type $B_m$ or of type $F_4$, we confirm that the requirements to apply   theories in Section 2 and Section 3 hold. In Section \ref{sect:H2 cohomology}, we  compute the Lie algebra cohomology $H^2(\frak g_-, \frak g)$.
In Section \ref{sect:local equivalence}, we complete the proof of main Theorem.\\


\noindent{ \bf Acknowledgements.}
%
We would like to thank  Qifeng Li and Jun-Muk Hwang for  valuable discussions, particularly for sharing the ideas in the preprint for the   $G_2$-type case.
The first named author was supported by the Institute of Basic Science IBS-R032-D1. The second named author was  supported by the Institute of Basic Science IBS-R003-D1.


\section{Varieties of minimal rational tangents} \label{sect:vmrt}

\subsection{Definitions and properties}

We recall definitions and properties of varieties of minimal rational tangents. For details, see  Section 1 of \cite{HM99} and Section 1 of \cite{Mk08}.

\begin{definition} Let $M$ be a uniruled projective manifold and let $\mathcal K$ be a minimal rational component, i.e., an irreducible component of the space of rational curves on $M$ such that the rational curves in $\mathcal K$ sweep out a Zariski open subset of $M$ and, with respect to a fixed ample line bundle on $M$, the degree of the rational curves in $\mathcal K$  is minimal among all such irreducible components. Any rational curve $C$ in $\mathcal K$ can be represented by a parameterized rational curve $f: \mathbb P^1 \rightarrow M$ so that $C=f(\mathbb P^1)$.

 Denote    by $\rho:\mathcal U \rightarrow \mathcal K$ and $\mu:\mathcal U \rightarrow M$  the universal family associated to $\mathcal K$.
 We call a rational curve $C$ \emph{standard} if $f^*TM$ is decomposed as  $f^*TM=\mathcal O(2) \oplus \mathcal O(1)^p \oplus \mathcal O^q$ where $f:\mathbb P^1 \rightarrow M$ is a parameterized rational curve representing $C $.
Then there is a smallest closed subvariety $E \subset M$   such that for every $x \in M\backslash E$ generic rational curve $C$ in $\mathcal K$ passing through $x$ is standard, i.e.,
We call $E$    the \emph{bad locus} of $\mathcal K$

Define the tangent map $\tau: \mathcal U \dashrightarrow \mathbb P(TM)$ by $\tau([f])=[df(T_0\mathbb P^1)]$ for $f:\mathbb P^1 \rightarrow M$ with $df(0)\not=0$.
For $x \in M$ the image $\mathcal C_x$ of the rational map $\tau_x: \mathcal U_x:=\mu^{-1}(x) \dashrightarrow \mathbb P(T_xM)$ is called the \emph{variety of minimal rational tangents} of $M$ at $x$.
The proper image $\mathcal C \subset \mathbb PTM$  of $\tau: \mathcal U \dashrightarrow \mathbb P TM$ is called the
(compactified) total space of varieties of minimal rational tangents.

\end{definition}

We will use the following results on varieties on minimal rational tangents (Proposition \ref{general mrc not intersect  codimension 2 subvariety}  -- \ref{Cartan-Fubini}).

\begin{proposition}[Proposition II.3.7 of \cite{Kol}] \label{general mrc not intersect  codimension 2 subvariety} Let $M$ be a uniruled projective manifold and $\mathcal K$ be a minimal rational component on $M$. Let $Z \subset M$ be a subvariety of codimension $\geq 2$. Then   a general $C \in \mathcal K$   does not intersect $Z$.
\end{proposition}


\begin{proposition} [Proposition 13 of \cite{HM98}, Proposition 1.2.2 of \cite{HM99}] \label{non-integrability}   Let $M$ be a uniruled projective manifold of Picard number one and $\mathcal K$ be a minimal rational component on $M$.
Assume that the variety $\mathcal C_x$ of minimal rational tangents at a general point $x \in M$ is linearly degenerate. Let $D$ be a proper meromorphic distribution on $M$ which contains the linear span of $\widehat{\mathcal C}_x$ at a general point. Then $D$ cannot be integrable.
\end{proposition}


\begin{theorem} [Theorem 1.2 of \cite{HM01}] \label{Cartan-Fubini} Let $M_1$ and $M_2$ be  uniruled projective manifolds of Picard number one and $\mathcal K_1$ and $\mathcal K_2$ be  minimal rational component on $M_1$ and $M_2$, respectively.
Let $E_1$ and $E_2$ be the bad locus of $\mathcal K_1$ and $\mathcal K_2$ and $\mathcal C_1 $ and $\mathcal C_2$ be varieties of minimal rational tangents of $(M_1, \mathcal K_1)$ and $(M_2, \mathcal K_2)$.
Assume that the general fiber $\mathcal C_x$  is positive dimensional and the Gauss map on
$\mathcal C_x$ is generically finite.

Let $U_1 \subset M_1 \backslash E_1$ and $U_2 \subset M_2 \backslash E_2$ be connected open subset and $f: U_1\rightarrow  U_2$ be a biholomorphism such that $[df](\mathcal C_1|_{U_1}) = \mathcal C_2|_{U_2}$.
Then there is a unique biholomorphic map $F: M_1 \rightarrow M_2$ such that $F|_{U_1} = f$.

\end{theorem}



\subsection{Fundamental forms}
Fundamental forms are basic invariants of projective varieties. We investigate behaviors of relative second and third fundamental forms of the varieties of minimal rational tangents along the liftings of standard minimal rational curves.
\begin{definition}
Let $U$ be a vector space and $Z \subset \mathbb P(U)$ be a subvariety. For a point $z$ in the smooth locus $Z^0$ of $Z$, denote by $\widehat T_zZ$ the affine tangent space of $Z$ at $z$. Then the intrinsic tangent space  at $z$ is $T_zZ=\widehat z ^* \otimes  \widehat T_zZ/\widehat z$ and the normal space is $N_z = T_z(\mathbb P(U))/T_zZ = \widehat z^* \otimes U/\widehat T_zZ$, where $\widehat z$ is the one-dimensional subspace of $U$ corresponding to the point $z$. The Gauss map $\gamma: Z^0 \rightarrow Gr(n+1, U)$ is defined by sending $z \in Z^0$ to the affine tangent space $\widehat T_zZ$ of $Z$ at $z$.  The differential $d \gamma$ defines a linear map $II_z:S^2T_zZ \rightarrow N_z$, called the {\it second fundamental form} of $Z$ at $z$. The image of $II_z$ is called the {\it second normal space} $N^{(2)}_z$ of $Z$ at $z$. $\widehat z \otimes Im\, II_z$ is a subspace of $U/\widehat T_zZ$ and the {\it second osculating affine tangent space} $T^{(2)}_zZ$ is defined by the subspace of $U$ whose quotient space by $\widehat T_zZ$ is   $\widehat z \otimes  Im \,II_z $.

Let $Z^{00}$ denote the points in $Z$ where the rank of the second fundamental form does not drop.
  The second Gauss map $\gamma^{(2)}: Z^{00} \rightarrow Gr(n^{(2)}+1, U)$ is defined by sending $z \in Z^{00}$ to the second osculating affine tangent space $\widehat T^{(2)}_zZ$ of $Z$ at $z$.  The differential $d \gamma^{(2)}$ defines a linear map $III_z:S^3T_zZ \rightarrow \widehat x^* \otimes U/T^{(2)}_zZ = N_z/N^{(2)}_z$, called the {\it third fundamental form}  of $Z$ at $z$. For any $k \geq 4$ the $k$-th fundamental form is defined in a similar way.
\end{definition}

\begin{definition}

Let $\mathcal U$ be a   vector bundle on a manifold $M$ and $\pi: \mathbb P \mathcal U \rightarrow M$ be the projection map from its projectivization. Let $\mathcal Z \subset \mathbb P(\mathcal U)$ be a subvariety  and let $\varpi: \mathcal Z \rightarrow   M$ be the restriction of $\pi$ to $\mathcal Z$.

Let $\mathcal Z^0\subset \mathcal Z$ be an open subset  such that $\varpi|_{\mathcal Z^0}: \mathcal Z^0 \rightarrow M^0:=\varpi(\mathcal Z^0)$ is a submersion and for each $t \in M^0$, $\mathcal Z_t:=   \varpi^{-1}(t) $ is immersed in $\mathbb P \mathcal U_t:=\pi^{-1}(t)$ at any point in $ \mathcal Z^0_t:= \varpi^{-1}(t) \cap \mathcal Z^0$.
Define a vector bundle $T^{\varpi}$ on $\mathcal Z^0$ by $T^{\varpi}:=\cup_{z\in \mathcal Z^0} T_{z} \mathcal Z_{\varpi(z)}$, which is called the \emph{relative tangent bundle} of $\mathcal Z \subset \mathbb P(\mathcal U)$. In a similar way we can define the relative affine tangent bundle $\widehat {T}^{\varpi}$,  relative normal bundles $\mathcal N $ and  relative second  fundamental form  $II^{\varpi}$. Then
$II^{\varpi}$ is a section of $\Hom(\Sym^2 T^{\varpi}, \mathcal N)$.

Assume that the rank of $II^{\varpi}$ is constant. Then the image of $II^{\varpi}$ defines a subbundle $\mathcal N^{(2)}$ of $\mathcal N$, called the \emph{second  normal bundle}. Similarly, we can define the relative third fundamental form $III^{\varpi}$ as a section of $\Hom(\Sym^3 T^{\varpi}, \mathcal N/\mathcal N^{(2)})$.

\end{definition}


Let $M$ be a uniruled projective manifold and $ \mathcal C$ be the variety of minimal rational tangents associated to a choice of a minimal rational component $\mathcal K$ on $M$.
For a rational curve $C$ in $\mathcal K$ represented by $f: \mathbb P^1 \rightarrow M$, we will denote by  $f^{\sharp}$   the map $\mathbb P^1 \rightarrow \mathbb P(TM)$ mapping $z \in \mathbb P^1$ to $df(T_z\mathbb P^1)$ and by $C^{\sharp}$ the image $f^{\sharp}(\mathbb P^1)$ of $f^{\sharp}$.

\begin{proposition} [Proposition 2.2 of \cite{Mk08}] \label{parallel transport of II}

Let $M$ be a uniruled projective manifold and let  $ \mathcal C$ be the variety of minimal rational tangents associated to a choice of a minimal rational component $\mathcal K$ on $M$. Denote by $\varpi: \mathcal C \rightarrow M$   the restriction of the projection map $\mathbb P(TM) \rightarrow M$.
Let $C=[f]$ be a standard rational curve on $X$ so that $f^*TM =\mathcal O(2) \oplus \mathcal O(1)^p \oplus \mathcal O^q$ for some $p$ and $q$.
Then
\begin{enumerate}
\item the pull-back $(f^{\sharp})^*\widehat{T}^{\varpi} $ of the relative affine tangent bundle $ \widehat{T}^{\varpi} $ of $\mathcal C \subset \mathbb P(TM)$ is the positive part $ P:=\mathcal O(2) \oplus \mathcal O(1)^p$ of $f^*TM$;
\item the relative 2nd fundamental form $II^{\varpi}$ of $\mathcal C|_{C^{\sharp}}$ is constant and  the pull-back $(f^{\sharp})^*\widehat{T}^{(2), \varpi} $ of the relative  second osculating affine   bundle   $ \widehat{T}^{(2), \varpi} $ of $\mathcal C \subset \mathbb P(TM)$  is a subbundle  $P^{(2)} =\mathcal O(2) \oplus \mathcal O(1)^p \oplus \mathcal O^r$ of $f^*TM$, where $r$ is the dimension of the image of $II^{\varpi}$.

\end{enumerate}

\end{proposition}

Proposition \ref{parallel transport of II} implies that the second fundamental forms of the varieties of minimal rational tangents are constant along the lifting of standard minimal rational curves. However, the third fundamental forms can vary. We will show the consistency of the third fundamental forms under some assumptions  (Proposition \ref{parallel transport of III}).

Let $D$ be the distribution on $M$ defined by the linear span of the affine cone $\widehat{\mathcal C}_x$.  Then outside a subvariety $Sing(D)$ of codimension $\geq 2$, $D$ is a subbundle of $TM$ and we may think of relative fundamental forms of $\mathcal C \subset \mathbb P(D)$.
   The kernel of the Frobenius bracket $[\,,\,]: D \wedge  D \rightarrow TM/D$
   is related to the projective geometry of the variety of minimal rational tangents as follows.
\begin{proposition} [Proof of Proposition 1.2.1 and  Proposition 1.3.1 and Proposition 1.3.2 of \cite{HM99}] \label{kernel of frobenius braket}
 Let $M$ be a uniruled projective manifold and $\mathcal K$ be a minimal rational component on $M$. Assume that at a general point $x \in M$ the variety of minimal rational tangents $\mathcal C_x$ is linearly degenerate. Denote by $D_x \subset T_xM$ the linear span of the affine cone $\widehat{\mathcal C}_x$ of $\mathcal C_x$ and by  $[ \,,\,]: \wedge^2 D_x \rightarrow T_xM/D_x$ the Frobenius Lie bracket. Then
 for a generic $\alpha \in \widehat{\mathcal C}_x$ and $\xi , \eta \in T_{\alpha}\widehat{\mathcal C}_x$,
 \begin{eqnarray*}  \alpha \wedge \xi, \,\, \alpha \wedge II(\xi, \eta),\,\,\xi \wedge \eta
  \end{eqnarray*} are contained in the kernel of $[ \,,\,]: \wedge^2 D_x \rightarrow T_xM/D_x$.
 \end{proposition}

Proposition \ref{kernel of frobenius braket} implies that the second osculation affine tangent space $ \widehat  T^{(2)} _{\alpha}\mathcal C_x$ is contained in the kernel of $[\alpha, \,]: D_x \rightarrow T_xM/D_x$, and thus the rank of $[\alpha, \,]: D_x \rightarrow T_xM/D_x$ is less than or equal to the codimension of $ \widehat  T^{(2)} _{\alpha}\mathcal C_x$ in $D_x$.


\begin{proposition} [cf. Proposition 3.1 of \cite{Mk08}, Proposition 5.1 of \cite{HH}]  \label{parallel transport of III}
Let $M$ be a uniruled projective manifold of Picard number one and let  $ \mathcal C$ be the variety of minimal rational tangents associated to a choice of a minimal rational component $\mathcal K$ on $M$.
 Let $D$ be the distribution on $M$ defined by the linear span of the affine cone $\widehat{\mathcal C}_x$.
By Proposition \ref{general mrc not intersect  codimension 2 subvariety} we may take a standard minimal rational curve $C=f(\mathbb P^1)$ with $C \cap Sing(D) =\emptyset$.
 Let $s$ be the codimension of the second osculating affine tangent space of $\mathcal C_x$ in $\mathbb P(D_x)$.
 Assume that
  \begin{enumerate}
   \item the third fundamental form of $\mathcal C_x \subset \mathbb P(D_x)$ is surjective for generic $x \in C$;

   \item the rank of $[\mathcal O(2)_x, \,]: D_x \rightarrow T_xM/D_x$ is $s$ for generic $x \in C$.
\end{enumerate}

Then the relative 3rd fundamental form  $III^{\varpi}$ is constant along $C^{\sharp}$
and we have  $f^*D =\mathcal O(2) \oplus \mathcal O(1)^p \oplus \mathcal O^r \oplus \mathcal O(-1)^s$ and $f^*TM/D =\mathcal O(1)^s \oplus \mathcal O^t$.

\end{proposition}


\begin{proof}   If $s =0$, then $D$ is integrable and thus $D=TM$  (Proposition \ref{non-integrability}).
Assume that $s >0$.
Since $\widehat {\mathcal C}_x $ is contained in $D_x$, $P^{(2)}$ is a subbundle of $f^*D$. Furthermore, $f^*D/P^{(2)}$ is a subbundle of $f^*TM/P^{(2)}\simeq \mathcal O^{q-r}$ and thus is $\mathcal O(b_1)\oplus \dots \mathcal O(b_s) ,$  where $0 \geq b_1 \geq\dots \geq b_s$. It follows that the short exact sequence $0 \rightarrow P^{(2)}\rightarrow f^*D \rightarrow f^* D/P^{(2)} \rightarrow 0$ is split and
$$f^*D = \mathcal O(2) \oplus \mathcal O(1)^p \oplus  \mathcal O^r \oplus \mathcal O(b_1)\oplus \dots \mathcal O(b_s) ,$$ where $0 \geq b_1 \geq\dots \geq b_s$.

If all $b_j$ are zero, then the Frobenius Lie bracket $[\,,\,]: D \times D \rightarrow TM/D$ is zero (for,...) and thus $D$ is integrable, a contradiction. Therefore, $b_s<0$.

The relative 3rd fundamental form $III^{\varpi}$ of $\mathcal C|_{C^{\sharp}}$ is a section of $$\Hom(\Sym^3 T^{\varpi}, Hom(f^*L, f^*D/P^{(2)}))|_{\mathcal C^{\sharp}} \subset \Hom(\Sym^3 T^{\varpi}, N/N^{(2)})|_{\mathcal C^{\sharp}}, $$
  which is isomorphic to $\Hom(\Sym^3\mathcal O(-1)^p,  \mathcal O(b_1-2) \oplus \dots \oplus \mathcal O(b_s -2))$. By the surjectivity of the third fundamental forms $III$, we have $b_s \geq -1$. Thus $b_s=-1$ and $f^*D = \mathcal O(2) \oplus \mathcal O(1)^p \oplus \mathcal O^{r'}  \oplus \mathcal O(-1)^{s'}$ for some $r' \geq r$ and $s' \leq s$.

The Chern number of $f^*(TM/D)$ is $s'$. Since every factor of $f^*(TM/D)$ has nonnegative degree, the rank of the positive part $f^*(TM/D)_+$ of $f^*(TM/D)$ is $\leq s'$.

On the other hand, under the Frobenius bracket $[\,,\,] :D \times D \rightarrow TM/D$, the image $[\mathcal O(2), f^*D ] $ is contained in $f^*(TM/D)_+$ and has dimension $\leq s' \leq s$. By the condition that the rank of $[\mathcal O(2)_x, \,]:D_x \rightarrow T_xM/D_x$ is $s$, we have  $s \leq s'$ and thus we have $s=s'$. Consequently, $r=r'$ and $f^* D= \mathcal O(2) \oplus \mathcal O(1)^p \oplus \mathcal O^r \oplus \mathcal O(-1)^s$ and $f^*(TM/D) = \mathcal O(1)^s \oplus \mathcal O^{q-r-2s}$. Therefore, $III^{\varpi}$ is a section of a trivial vector bundle and thus is constant.
\end{proof}




\section{Geometric structures} \label{sect:geometric structures}

\subsection{$G_0$-structures and ${\bf S}$-structures}

\begin{definition}\label{distribution type} Let $\mathfrak m=\bigoplus_{p<0}\mathfrak g_{p}$ be a fundamental graded Lie algebra. A \emph{distribution} $D$ on a complex manifold $M$ is a subbundle of the tangent bundle $TM$ of $M$. A distribution $D$ is called {\it of type $\mathfrak m$} if for each $x \in M$ the symbol algebra $\Symb_x(D)$ is isomorphic to $\mathfrak m$ as a graded Lie algebra. In this case, the pair $(M,D)$ is called a \emph{filtered manifold of type} $\mathfrak m$.
\end{definition}


\begin{definition}  \label{G_0 structure} Keep notations as in Definition \ref{distribution type}.
For each $x \in M$, let $\mathscr R_x $ be the set of all isomorphisms $r \colon \mathfrak m\rightarrow \Symb_x(D)$ of graded Lie algebras. Then $\mathscr R:=\cup_{x\in M}\mathscr R_x $ is a principal $G_0(\mathfrak m)$-bundle on $M$, where  $G_{0}(\mathfrak m)$ is the automorphism group  of   the graded Lie algebra $\mathfrak m$.
We call   $\mathscr R$  the \emph{frame bundle} of $(M,D)$.

Given a closed subgroup $G_0\subset G_{0}(\mathfrak m)$, a \emph{$G_0$-structure} on $(M,D)$ is a $G_0$-subbundle of the frame bundle $\mathscr R$ of $(M,D)$. Two $G_0$-structures $\mathscr P_1$ on $(M_1,D_1)$ and $\mathscr P_2$ on $(M_2,D_2)$ are \emph{equivalent} if there is a biholomorphism $\varphi: M_1 \rightarrow M_2$ such that $d\varphi: TM_1 \rightarrow TM_2$ induces an isomorphism from $\mathscr P_1$ onto $\mathscr P_2$.
The local equivalence of two $G_0$-structures is defined similarly for open sets $U_1 \subset M_1$ and $U_2 \subset M_2$.
 \end{definition}

\begin{definition}\label{cone structure}  Keep notations as in Definition \ref{distribution type}. Let ${\bf S}$ be a nondegenerate subvariety of $\mathbb P {\frak g_{-1}} $.
A fiber subbundle $\mathcal S \subset \mathbb P D$ is called an  {\it ${{\bf S}}$-structure on $(M,D)$} if for each $x \in M$, the fiber $\mathcal S_x \subset \mathbb P D_x$ is isomorphic to ${\bf S} \subset \mathbb P {\frak g_{-1}}$ under a graded Lie algebra isomorphism
$  \mathfrak m \rightarrow   \Symb_x(D) $.

Two ${\bf S}$-structures $\mathcal S_1$ on $(M_1,D_1)$ and $\mathcal S_2$ on $(M_2,D_2)$ are said to be \emph{equivalent} if there exists a biholomorphism $\phi \colon M_1 \to M_2$ such that $d\phi \colon \mathbb P TM_1 \to \mathbb P TM_2$ sends $\mathcal S_1 \subset \mathbb P TM_1$ to $\mathcal S_2 \subset \mathbb P TM_2$. The \emph{local equivalence} of two ${\bf S}$-structures is defined similarly for   open subsets $U_1\subset M_1$ and $U_2 \subset M_2$.
\end{definition}

\begin{definition}   Let ${\bf U}$ be a vector space and and let $ {\bf S} $ be a nondegenerate subvariety of $\mathbb P \bf U$. Consider the graded free Lie algebra $F({\bf U})$ generated by $\bf U$ such that  $F({\bf U})_{i}:=\bf U ^{(i)}=[\cdots[[\bf U,\bf U],\bf U],\bf U \cdots]$, and  the ideal $I$ of $F(\bf U)$ generated by the relation $[v,w]=0$ for $v,w \in {\bf U}$ such that $v \in \widehat{{\bf S}}$ and  $w \in T_v \widehat{{\bf S}}$. Keep notations as in Definition \ref{cone structure}, i.e. let $\mathfrak m=\bigoplus_{p<0}\mathfrak g_{p}$ be a fundamental graded Lie algebra and let ${\bf S}$ be a nondegenerate subvariety of $\mathbb P {\frak g_{-1}}$. If the  quotient graded Lie algebra $ F({\frak  g_{-1}})/I$ is isomorphic to $\frak m$, we say that the fundamental graded Lie algebra $\frak m$ is  \emph{determined by} ${\bf S} \subset \mathbb P \frak g_{-1}$.
\end{definition}

Let $G(\widehat{\bf S})$ be {\it the linear automorphism group of $\widehat{\bf S} \subset \frak g_{-1}$}, i.e., the  subgroup of $\GL(\frak g_{-1})$ consisting of   linear automorphism  of $\widehat{\bf S} \subset \frak g_{-1}$. Then $G(\widehat{\bf S})$ acts on $F({\frak  g_{-1}})/I$ preserving the graded Lie algebra structure. If $\mathfrak m=\bigoplus_{p<0}\mathfrak g_{p}$ is a fundamental graded Lie algebra determined by  $ {\bf S} $, this $G(\widehat{\bf S})$ action defines a homomorphism $G(\widehat{\bf S}) \rightarrow G_0(\frak m)$ by the isomorphism between $F({\frak  g_{-1}})/I$ and $\frak m$. The induced map $G(\widehat{\bf S}) \rightarrow G_0(\frak m)$ is injective and let $G_0 \subset G_0(\frak m)$ denote its image. Then the relation between $G_0$-structures and ${\bf S}$-structures is given as follows.


\begin{proposition} \label{S-str and G_0-str} Let $\mathfrak m=\bigoplus_{p<0}\mathfrak g_{p}$ be a fundamental graded Lie algebra determined by a nondegenerate subvariety  $ {\bf S} $    of $\mathbb P {\mathfrak g_{-1}}$. Let $G(\widehat{\bf S})$ be \it the linear automorphism group of $\widehat{\bf S} \subset \frak g_{-1}$ and consider the induced map $G(\widehat{\bf S}) \rightarrow G_0(\frak m)$ and let $G_0 \subset G_0(\frak m)$ denote its image.
 Then, there is a one-to-one correspondence between $G_0$-structures  and ${\bf S}$-structures on filtered manifolds of type $\mathfrak m$. Furthermore, two $G_0$-structures $\mathscr P_1$ and $\mathscr P_2$ are equivalent if and only if the corresponding ${\bf S}$-structures $\mathcal S_1$ and $\mathcal S_2$ are equivalent.
\end{proposition}

\begin{proof} Let $(M,D)$ be a filtered manifold of type $\mathfrak m$.
Given   a $G_0$-structure  $\mathscr P$  on $(M,D)$, define $\mathcal S_x$ by $[r] ({\bf S})$ for any $r \in \mathscr P_{x}$, where $[r]$ is the isomorphism $\mathbb P {\frak g_{-1}}  \rightarrow \mathbb P D_x$ induced by the isomorphism  $r:\mathfrak m \rightarrow \Symb_x(D)$. Then $\mathcal S_x$ is well defined and the union $\mathcal S= \cup_{x \in M} \mathcal S_x$ defines an ${\bf S}$-structure on $(M,D)$.

Conversely, let $\mathcal S \subset \mathbb P D$ be an ${\bf S}$-structure on $(M,D)$. Define $\mathscr P_x$ by $\{ r \in \mathscr R_x: \varphi ({\bf S}) = \mathcal S_x \}$. Then $\mathscr P_x$ is well defined and the union $\mathscr P= \cup_{x \in M} \mathscr P_x$ defines a $G_0$-subbundle of the frame bundle $\mathscr R$ of $(M,D)$.
\end{proof}


\subsection{Lie algebra cohomologies}

\begin{notation} Let $\frak m=\bigoplus_{j <0} \frak g_j$ be a nilpotent Lie algebra and $\Gamma$ be a representation space of $\frak m$.
Define a complex
\begin{eqnarray*}
0 \stackrel{\partial}{\longrightarrow} \Gamma \stackrel{\partial}{\longrightarrow} \Hom( \frak m, \Gamma) \stackrel{\partial}{\longrightarrow} \Hom(\wedge^2 \frak m  ,\Gamma  ) \stackrel{\partial}{\longrightarrow} \dots
\end{eqnarray*}
by
\begin{eqnarray*}
\partial \phi (X_1, \dots,  X_{q+1}) &=& \sum_{i=1}^{q+1} (-1)^{i+1}X_i.\phi(X_1, \dots,\hat X_i, \dots , X_{q+1}) \\
&& + \sum_{1 \leq i<j \leq q+1} (-1)^{i+j} \phi([X_i, X_j], X_1, \dots, \hat X_i, \dots, \hat X_j , \dots , X_{q+1})
\end{eqnarray*}
 for $\phi \in \Hom(\wedge^q \frak m, \Gamma)$ and $X_1, \dots, X_{q+1} \in \frak m$.
The cohomology space
$$H^{q}(\frak m, \Gamma):=\frac{\text{Ker}\, (\partial : \Hom(\wedge^q \frak m  , \Gamma) \rightarrow \Hom(\wedge^{q+1}\frak m,   \Gamma))}
{\text{Im}\, (\partial : \Hom( \wedge^{q-1}\frak m ,\Gamma) \rightarrow \Hom(\wedge^q\frak m, \Gamma))}
$$ is called the {\it Lie algebra cohomology space associated to the representation of} $\frak m$ on $\Gamma$.

Assume that $\frak m$ and $\Gamma$ have  gradations such that $\frak m_i.\Gamma_p \subset \Gamma_{i+p}$.  Then $H^q(\frak m, \Gamma)$ also has a gradation
$$H^q(\frak m, \Gamma) = \oplus_{p } H^{p,q}(\frak m, \Gamma)$$
where
$$\wedge^q_j\frak m =\sum_{j_1+\dots+j_q=j, j_1,\dots, j_q <0} \frak g_{j_1} \wedge \dots \wedge \frak g_{j_q}$$
and
\begin{eqnarray*}
H^{p,q}(\frak m, \Gamma) = \frac{\text{Ker}\,(\partial :  \Hom(\wedge^q_j \frak m , \Gamma_{j+p+q-1}) \rightarrow \Hom(\wedge^{q+1}_j \frak m, \Gamma_{j+p+q-1}))}
{\text{Im}\,  (\partial :  \Hom(\wedge^{q-1}_j \frak m, \Gamma_{j+p+q-1}) \rightarrow \Hom(\wedge^{q }_j \frak m, \Gamma_{j+p+q-1})}.
\end{eqnarray*}
We also denote $H^{p,q}(\frak m, \Gamma)$ by $H^q(\frak m,  \Gamma)_{p+q-1}$.
\end{notation}


 \subsection{Prolongation methods} \label{sect:prolongation methods}

 We review the theory of prolongation methods in \cite{HM}. Let $\frak m=\bigoplus_{p<0} \frak g_p$ be a fundamental graded Lie algebra and $G_0$ be a connected subgroup of $G_0(\frak m)$ with Lie algebra $\frak g_0$. Let $\frak g=\bigoplus_{\ell \in \mathbb Z} \frak g_{\ell}$ be the prolongation of $(\frak m, \frak g_0)$.
 Then,    $\frak g_{\ell+1}$ ($\ell \geq 0$) can be regarded as a subspace of $F^{\ell+1}\frak{gl}(E^{(\ell)})/F^{\ell+2}$,  where $E^{(\ell)}:=\oplus_{i \leq \ell}\frak g_i$, %and is the Lie algebra of a Lie group $G_{\ell+1}$
 % Then
  and    corresponds to an abelian    group  $G_{\ell+1}$ ($\ell \geq 0$)  which can be regarded as a subgroup of $ F^{\ell+1}GL(E^{(\ell)})/F^{\ell+2}$.

  We call a sequence $$\mathscr P^{(\ell)} \longrightarrow \mathscr P^{(\ell-1)} \longrightarrow \dots \longrightarrow \mathscr P^{(0)} \longrightarrow M$$ of principal bundles   a {\it geometric structure of order} $\ell$ {\it of type} $(\frak g_-, G_0, \dots, G_k)$ if   $\mathscr P^{(i+1)} \rightarrow \mathscr P^{(i)}$ is a principal $G_{i+1}$-subbundle of the universal  frame bundle $\mathscr S^{(i+1)} \mathscr P^{(i)} \rightarrow \mathscr P^{(i)}$  of $\mathscr P^{(i)}$ of order $i+1$, where $0 \leq i \leq \ell$.
   For the definition of the universal frame bundle $\mathscr S^{(i+1)} \mathscr P^{(i)}$ of $\mathscr P^{(i)}$ of order $i+1$, see Definition 2.1 of  \cite{HM}.  The property we use is that a map between two geometric structures $\mathscr P^{(i)} $, $\mathscr Q^{(i)}$ of order $i$ induces a map between their universal frame bundles $\mathscr S^{(i+1)}\mathscr P^{(i)}$, $ \mathscr S^{(i+1)}\mathscr Q^{(i)}$ of order $i+1$.

 We write a  sequence $\mathscr P^{(\ell)} \longrightarrow \mathscr P^{(\ell-1)} \longrightarrow \dots \longrightarrow \mathscr P^{(0)} \longrightarrow M$ of principal bundles simply as $\mathscr P^{(\ell)}$ when no confusion can arise, and define the equivalence of two geometric structures $\mathscr P^{(\ell)} $ and $\mathscr Q^{(\ell)} $  inductively as follows.
  %
Two geometric structures $\mathscr P^{(\ell)}$ and $\mathscr Q^{(\ell)}$ are {\it equivalent} if their truncations $\mathscr P^{(\ell-1)}$ and $\mathscr Q^{(\ell-1)}$ are equivalent and the lifting $\mathscr S^{(\ell)} \mathscr P^{(\ell-1)} \rightarrow \mathscr S^{(\ell)} \mathscr Q^{(\ell-1)}$ of the equivalence $\mathscr P^{(\ell-1)} \rightarrow \mathscr Q^{(\ell-1)}$  maps $\mathscr P^{(\ell)}$ onto $\mathscr Q^{(\ell)}$. \\

Fix a set of subspaces $W=\{W^1_{\ell}, W^2_{\ell+1}\}_{\ell \geq 0}$ such that
\begin{eqnarray*}
 \Hom (\frak g_-,  \frak g )_{\ell} &= &W_{\ell}^1 \oplus \partial  \frak g_{\ell}\\
 \Hom(\wedge^2\frak g_-,   \frak g )_{\ell+1} &= &W_{\ell+1}^2 \oplus \partial \Hom(\frak g_-,  \frak g )_{\ell+1} .
 \end{eqnarray*}


\begin{theorem} [Theorem 3.1  of \cite{HM}] \label{thm:prolongation methods}

  Let $\mathscr P$ be a $G_0$-structure on a filtered manifold $(M,D)$ of type $\frak m$.
  Then for each $\ell \geq 1$, there is a geometric structure
 $$  \mathscr S_W^{(\ell)}\mathscr P \stackrel{G_{\ell}}{\longrightarrow} \mathscr S_W^{(\ell-1)}\mathscr P \longrightarrow \dots \longrightarrow \mathscr S_W^{(1)}\mathscr P \stackrel{G_1}{\longrightarrow} \mathscr P \stackrel{G_0}{\longrightarrow} M $$
 of type $(\frak g_-, G_0, \cdots, G_{\ell})$.
 %such that each $\mathscr S_W^{(i+1)}\mathscr P \stackrel{G_{i+1}}{\longrightarrow} \mathscr S_W^{(i)}\mathscr P$ is a principal subbundle of the frame bundle of $\mathscr S_W^{(i)}\mathscr P$ of order $i+1$, where $0 \leq i \leq \ell$.
  %
  Furthermore, two $G_0$-structures $\mathscr P$ and $\mathscr Q$ are   equivalent if and only if the corresponding geometric structures $\mathscr S_W^{(\ell)}\mathscr P$ and $\mathscr S_W^{(\ell)}\mathscr Q$ are  equivalent.
\end{theorem}

We call the limit $\mathscr S_W \mathscr P=\lim_{\ell} \mathscr S_W^{(\ell)}\mathscr P$ the {\it $W$-normal complete step prolongation of} $\mathscr P$. \\

Now
 assume that  $\frak g$ is finite dimensional. Then there is a Lie group $G$ and its subgroup $G^0$ which contains $G_0$, with Lie algebras $\frak g$ and $\frak g^0:=\bigoplus _{\ell \geq 0} \frak g_{\ell}$. %, such that $G \rightarrow G/G^0$ is a $G_0$-structure on a filtered manifold $G/G^0$ of type $\frak m$.
 By {\it a Cartan connection of type} $G/G^0$, we mean a principal $G^0$-bundle $P$ on a manifold $M$ with a $\frak g$-valued 1-form $\theta$ on $P$ satisfying the following properties.
\begin{enumerate}
\item $\theta: T_zP \rightarrow \frak g$ is an isomorphism for all $z  \in P$
\item $R_a^* \theta = \mathrm{Ad} (a)^{-1}\theta$ for $a \in G^0$
\item $\theta(\widetilde A) =A $ for $A \in \frak g^0$.
\end{enumerate}

Given a $G_0$-structure $\mathscr P$ on a filtered manifold  $(M,D)$  of type $\frak m$, for $k \geq 1$,  define a vector bundle $\mathcal H^2_k$ on $M$  by  $\mathcal H^2_k:=\mathscr P \times _{G_0} H^2(\frak m, \frak g)_k$. % called {\it ...}.

\begin{proposition} [Theorem 7.4 of \cite{HM}] \label{prop:prolongation methods} Let a $\mathscr P$ be a $G_0$-structure on a filtered manifold  $(M,D)$  of type $\frak m$.
  If  $H^0(M, \mathcal H^2_k)$ is zero for all $k \geq 1$, then the $W$-normal complete step  prolongation  $\mathscr S_W \mathscr P $ of $\mathscr P$ is a Cartan connection of type $G/G^0$ which is flat, and $\mathscr P$ is locally isomorphic to the standard $G_0$-structure on $G/G^0$.
\end{proposition}




\section{Smooth horospherical varieties of Picard number one} \label{horospherical varieties}

\subsection{Classifications}

Let $\mathfrak l$ be a semisimple Lie algebra. We fix a Cartan subalgebra $\mathfrak h$ of $\mathfrak l$ and let $\Phi$ be the set of roots of $\mathfrak l$ relative to $\mathfrak h$. The root space decomposition of $\mathfrak l$ is given by
 \begin{eqnarray*} \mathfrak l =\mathfrak h \oplus \bigoplus_{\alpha \in \Phi} \mathfrak l_{\alpha}, \end{eqnarray*}
where $\mathfrak l_{\alpha}
$ is the root space of $\alpha \in \Phi$.

\begin{definition}
Let $\{\alpha_1, \cdots, \alpha_m \}$ be a set of simple roots of $\mathfrak l$. We define the \emph{characteristic element  associated with $\alpha_i$} as an element $E_{\alpha_i}$ in $\mathfrak h$ such that $\alpha_j(E_{\alpha_i})= \delta_{i,j}$ for $i,j=1, \dots, m$. Define a gradation $\mathfrak l=\bigoplus_{p\in \mathbb Z} \mathfrak l_p$  on $\mathfrak l$ by $\mathfrak l_p = \{ v \in \mathfrak l : [E_{\alpha_i}, v] = p v\}$  for $p \in \mathbb Z$, which is called the \emph{gradation associated with $\alpha_i$}. In general, given a representation $V$ of $\frak l$ we define the \emph{gradation associated with $\alpha_i$} in a similar way.
\end{definition}

\begin{notation} Given a set $\{\alpha_1, \cdots, \alpha_m \}$ of simple roots of  $\frak l$, let $\{\varpi_1, \dots, \varpi_m\}$ be the set of fundamental weights. For each $i=1, \dots, m$, let $P^{\alpha_i}$ denote the maximal parabolic subgroup associated to $\alpha_i$ and let $V_{\varpi_i}$ denote the irreducible representation of Lie group $L$ with highest weight $\varpi_i$, where $L$ is semisimple Lie group corresponding to $\frak l$
\end{notation}


For a reductive   group $L$, a normal $L$-variety is said to be \emph{horospherical} if
it has an open $L$-orbit $L/H$ whose isotropy group $H$ contains the unipotent part of a Borel subgroup of $L$.


\begin{theorem}[Theorem 0.1 and Theorem 1.11 of \cite{Pa}]\label{Horospheical}
Let $L$ be a reductive group.
 Let $X$ be a smooth nonhomogeneous projective horospherical $L$-variety with Picard number one.
 Then $X$ is uniquely determined by its two closed $L$-orbits $ Y$ and $Z$, isomorphic to $L/P_Y$ and $L/P_Z$, respectively; and $(L, \alpha, \beta)$ in one of the triples of the following list, where $P_Y=P^\alpha$ and $P_Z=P^\beta$ for simple roots $\alpha$ and $\beta$.
 \begin{enumerate}
   \item[\rm(1)] $(B_m,{\alpha_{m-1}},{\alpha_{m}})$ for $m\geq3$;
   \item[\rm(2)] $(B_3,{\alpha_{1}},{\alpha_{3}})$;
   \item[\rm(3)] $(C_m,{\alpha_{i+1}},{\alpha_{i}})$ for $m\geq2$, $i\in \{1,\ldots, m-1\}$;
   \item[\rm(4)] $(F_4,{\alpha_{2}},{\alpha_3})$;
   \item[\rm(5)] $(G_2,{\alpha_{2}},{\alpha_1})$.
 \end{enumerate}
  Moreover, the automorphism group $\Aut(X)$ of $X$ is
  $(SO(2m+1)\times \mathbb C^*)\ltimes V_{\varpi_m}$,   $(SO(7)\times \mathbb C^*)\ltimes V_{\varpi_3}$,  $((Sp(2m)\times \mathbb C^*)/\{\pm1\})\ltimes V_{\varpi_1}$,    $(F_4\times \mathbb C^*)\ltimes V_{\varpi_4}$  and $(G_2\times \mathbb C^*)\ltimes V_{\varpi_1}$, respectively.

  Finally, $\Aut(X)$ has two orbits in $X$,   $Z$ and the complement of $Z$ in $X$.
\end{theorem}

By Theorem \ref{Horospheical}, $Y=L/P_Y$ is contained in the open orbit $X^0$ of $\Aut(X)$ in $X$. In particular,  the base point $o$ of $Y=L/P_Y$ is contained in $X^0$. We will take $o$ as the base point of the quasi-projective variety $X$.


\begin{proposition}[Section 2, Proposition 48 and Propositon 49 of \cite{Ki}]\label{tangentsp of X}
Let $X$ be a smooth nonhomogeneous projective horospherical variety $(L, \alpha, \beta)$ of Picard number one. Let $G=\Aut(X)$ and let $\mathfrak g=(\mathfrak l+\mathbb C) \rhd U$ be the corresponding Lie algebra.
Then,
 \begin{enumerate}
    \item[(1)] there is a grading on $\frak l$ and $U$,
\begin{eqnarray*}
\mathfrak l=\bigoplus_{k=-\mu }^{\mu}\mathfrak l_k \text{ \, and \, }  U=\bigoplus_{k=-1}^{\nu }U_k,
\end{eqnarray*}
such that the negative part $\mathfrak m=\bigoplus_{p<0}\mathfrak g_{p}$ of $\frak g$, where the   grading is defined    by
 \begin{eqnarray*}
 \mathfrak g_0&:=&(\mathfrak l_0\oplus \mathbb C)\rhd U_{ 0} \\
 \mathfrak g_p&:=&\mathfrak l_p \oplus U_{p} \text{ for } p \not=0,
 \end{eqnarray*}
 is identified with the tangent space of $X$ at the base point $o$ of $X$. Here, we assume that $\frak l_k=0$ for $k<-\mu$ or $k > \mu$, and $U_k=0$ for $k <-1$ or $k >\nu$.
 \item[(2)] $H^1(\frak m, \frak g)_{p}=0$ for $p>0$.
 \item[(3)] $\frak g=\bigoplus_{\ell \in \mathbb Z} \frak g_{\ell}$ is the prolongation of $(\frak m, \frak g_0)$.
 \end{enumerate}
\end{proposition}

Hence, we could consider the prolongation methods of section \ref{sect:prolongation methods} for a smooth nonhomogeneous projective horospherical variety of Picard number one

\subsection{Varieties of minimal rational tangents}

In   this section we will describe the varieties of minimal rational tangents of horospherical varieties in the list of Theorem \ref{Horospheical}. For an arbitrary reductive group $L$ and for a finitely many irreducible $L$-representation spaces $V_i$ ($i=1, \dots, r$),  let $\mathcal H_{L}(\oplus_{i=1}^r V_i)$ denote the closure of the sum of highest weight vectors $v_i$ of $V_i$ in $\mathbb P(\oplus_{i=1}^r V_i)$. For example, $\mathcal H_{L}(V)$ for an irreducible representation space $V$ is the highest weight orbit, and $(L, \alpha_i, \alpha_j)$ is $\mathcal H_{L}(V_{\varpi_i} \oplus V_{\varpi_j})$, where $V_{\varpi_i}$ (respectively, $V_{\varpi_j}$) is the irreducible representation of $L$ of highest weight $\varpi_i$ (respectively, $ \varpi_j$).

With the notions in Proposition \ref {tangentsp of X}, let $L_0$ be the subgroup of $G=\Aut(X)$ with Lie algebra $\mathfrak l_0$.

\begin{proposition} [Proposition 3.5.2 of \cite{Ki0}]\label{vmrt of X} Let $X$ be a smooth nonhomogeneous projective horospherical variety $(L, \alpha, \beta)$ of Picard number one.
 Let $\mathcal C_o(X) \subset \mathbb P (T_oX) $ denote  the variety of minimal rational tangents of $X$  at the base point $o$. Then
 \begin{eqnarray*}
 \mathcal C_o(X) = \left\{ \begin{array}{l}
 \mathcal H_{L_0}( U_{-1} \oplus\mathfrak l_{-1} \oplus \mathfrak l_{-2} ) \text{ if } X \text{ is }  (C_m, \alpha_{i+1}, \alpha_{i}) \text{ for }  1 \leq i < m \\
 \mathcal H_{L_0}(  U_{-1}\oplus \mathfrak l_{-1}), \text{ otherwise. }
 \end{array}\right.
   \end{eqnarray*}
\end{proposition}

\begin{proof} As we mentioned in previous subsection, we will take the base point $o$ of $Y=L/P_Y$ as the base point of $X$.  Then the variety $\mathcal C_o(Y)$ of minimal rational tangents of $Y$ at $o$  is contained in the variety $\mathcal C_o(X)$ of minimal rational tangents of $X$ at $o$.

\medskip
Assume that $X$ is not $(C_m, \alpha_{i+1}, \alpha_{i}), m>2, i=1,\ldots, m-2$. Since $Y=L/P_Y$ is associated with a long simple root $\alpha$, by Proposition 1 of Hwang-Mok (\cite{HM02}), the variety $\mathcal C_o(Y)$ of minimal rational tangents of $Y$ is  $\mathcal H_{L_0}(\mathfrak l_{-1})$. Furthermore, $L_0 \rhd U_{0}$ acts invariantly on $\mathcal C_0(X)$. From $[U_{0}, \frak l_{-1}] \subset U_{-1}$, it follows that the highest weight vector of $U_{-1}$ is contained in $\mathcal C_o(X)$. Therefore,  $\mathcal H_{L_0}( U_{-1}\oplus \mathfrak l_{-1})$ is contained in  $\mathcal C_o(X)$ so that $\dim \mathcal H_{L_0}( U_{-1}\oplus \mathfrak l_{-1})$ is less than or equal to $\dim \mathcal C_o(X)$. However, $\dim \mathcal C_o(X)$ cannot exceed $\dim H^0(C_{\alpha}, N_{C_{\alpha}|X} (-1))$ which is equal to   $K^{-1}_X.C_{\alpha}-2$. % of $X$.
Now by comparing the dimension of $\mathcal H_{L_0}( U_{-1}\oplus \mathfrak l_{-1})$ with   $K^{-1}_X.C_{\alpha} -2$, we get the desired results.
\begin{center}
    \begin{tabular}{ c c  c  }
types & $K^{-1}_X.C_{\alpha}-2$ \\
\hline
$(B_m, \alpha_{m-1}, \alpha_{m})$ & $m$ \\
$(B_3, \alpha_{1}, \alpha_{3})$ & 5 \\
$(C_m, \alpha_{m}, \alpha_{m-1})$ & $m$ \\
$(C_m, \alpha_{i+1}, \alpha_{i})_{ 1 < i+1 < m}$ & $2m-i-1$ \\
$(F_{4}, \alpha_{2}, \alpha_{3})$ & 4\\
$(G_{2}, \alpha_{2}, \alpha_{1})$ & 2 \\
    \end{tabular}
\end{center}
 Here, we use the following description of the representation $ U_{-1}\oplus \mathfrak l_{-1}$ of $L_0$ and $\mathcal H_{L_0}( U_{-1}\oplus \mathfrak l_{-1})$.

\begin{enumerate}

\item[(1)] $(B_m, \alpha_{m-1}, \alpha_{m})$, $m>2$ where $U=V_{\varpi_m}$; Denote $L_0=A_1\times A_{m-2}$. Let $V$ be the standard representation of $A_1$ and $W^*$ be the standard representation of $A_{m-2}$. Then
\begin{eqnarray*}
\mathfrak l_{-1}&=& \Sym^2 V \otimes W \\
U_{-1}&=& V.
\end{eqnarray*}
The closure  $\mathcal H_{L_0}( U_{-1}\oplus \mathfrak l_{-1})$ of the $L_0$-orbit at $v + v^2\otimes w$ is
\begin{eqnarray*}
 \mathcal H_{L_0}( U_{-1}\oplus \mathfrak l_{-1})&=&\mathbb P \{c v + v^2 \otimes w: c \in \mathbb C, v \in V, w \in W\} \\
 &\simeq& \mathbb P (\mathcal O_{\mathbb P(V)}(-1) \oplus \mathcal O_{\mathbb P(V)}(-2)^{m-1}),
\end{eqnarray*}
where $\dim V=2$ and $ \dim W=m-1$.

\item[(2)]$(B_3, \alpha_{1}, \alpha_{3})$ where $U=V_{\varpi_3}$; Denote $L_0=B_2$. Let $V$ be the spin representation of $B_2$. Let $W$ be the standard representation of $B_2$. Then
\begin{eqnarray*}
\mathfrak l_{-1}&=& W\\
U_{-1}&=& V.
\end{eqnarray*}
The closure $\mathcal H_{L_0}( U_{-1}\oplus \mathfrak l_{-1})$ of the $L_0$-orbit at $v+w$ is the horospherical variety of type $(C_2, \alpha_2, \alpha_1)$, the odd symplective Grassmannian $Gr_w(2,\mathbb C^5)$ of isotropic $2$-subspaces in $\mathbb C^5$.

\item[${\rm (3)_{m}}$] $(C_m, \alpha_{m}, \alpha_{m-1})$ where $U=V_{\varpi_1}$; Denote $L_0=A_{m-1}$. Let $V^*$ be the standard representation of $A_{m-1}$. %with $\dim V = m$.
Then
\begin{eqnarray*}
\mathfrak l_{-1}&=& \Sym^2V \\
U_{-1}&=& V.
\end{eqnarray*}
The closure  $\mathcal H_{L_0}( U_{-1}\oplus \mathfrak l_{-1})$ of the $L_0$-orbit at $v+ v^2$ is
 \begin{eqnarray*}
\mathcal H_{L_0}( U_{-1}\oplus \mathfrak l_{-1})&=& \mathbb P \{c v + v^2 : c \in \mathbb C, v \in V\} \\
&\simeq& \mathbb P (\mathcal O_{\mathbb P(V)}(-1) \oplus \mathcal O_{\mathbb P(V)}(-2)),
\end{eqnarray*}
where $\dim V=m$.

\item[(4)] $(F_{4}, \alpha_{2}, \alpha_{3})$ where $\alpha_{2}$ is a long root and $U=V_{\varpi_4}$; Denote $L_0=A_{1}\times A_2$. Let $V$ be the standard representation of $A_{1}$ and $V^*=V$ and let $W^*$ be the standard representation of $A_{2}$. Then
\begin{eqnarray*}
%\mathfrak l_{-3}&=& W \\
%\mathfrak l_{-2}&=& \Sym^2V \\
\mathfrak l_{-1}&=& \Sym^2V\otimes W \\
U_{-1}&=&V.
\end{eqnarray*}
%The number $K^{-1}_X.C_{\alpha} -2$   is $4$. % and is  equal to  the dimension of
The closure   $\mathcal H_{L_0}( U_{-1}\oplus \mathfrak l_{-1})$ of the $L_0$-orbit at $v+ v^2\otimes w$ is%, where  $v$ is a highest weight vector of $V^*$ and $w$ is a highest weight vector of $W^*$, is
%\begin{eqnarray*} \overline{L_0.(v^2\otimes w + v)} \subset \mathbb P ( \Sym^2V^*\otimes W^* + V^*)=\mathbb P(\mathfrak g_{-1}), \end{eqnarray*}
\begin{eqnarray*}
\mathcal H_{L_0}( U_{-1}\oplus \mathfrak l_{-1})&=& \mathbb P \{c v + v^2 \otimes w: c \in \mathbb C, v \in V, w \in W\} \\
&\simeq& \mathbb P (\mathcal O_{\mathbb P(V)}(-1) \oplus \mathcal O_{\mathbb P(V)}(-2)^{2}),
\end{eqnarray*}
where $\dim V=3$ and $ \dim W= 2$.

\item[(5)] $(G_{2}, \alpha_{2}, \alpha_{1})$ where $U=V_{\varpi_1}$; Denote $L_0=A_{1}$. Let $V$ be the standard representation of $A_{1}$ such that $V^*=V$. Then
\begin{eqnarray*}
%\mathfrak l_{-2}&=& \wedge^2V \\
\mathfrak l_{-1}&=& \Sym^3V \\
U_{-1}&=& V.
\end{eqnarray*}
%The number $K^{-1}_X.C_{\alpha} -2$  is $2$. % and is equal to the dimension of $L_0$-orbit %
The closure $\mathcal H_{L_0}( U_{-1}\oplus \mathfrak l_{-1})$ of the  $L_0$-orbit at $v+ v^3$ is  %, where $w$ is a highest weight vector of $W^*$, is
%\begin{eqnarray*} \overline{L_0.(w^3 + w)} \subset \mathbb P (\Sym^3 W^* + W^*)=\mathbb P(\mathfrak g_{-1}),\end{eqnarray*}
\begin{eqnarray*}
\mathcal H_{L_0}( U_{-1}\oplus \mathfrak l_{-1})&=& \mathbb P \{c v + v^3 : c \in \mathbb C, v \in V\} \\
&\simeq& \mathbb P (\mathcal O_{\mathbb P(V)}(-1) \oplus \mathcal O_{\mathbb P(V)}(-3)),
\end{eqnarray*}
where $\dim V=2$.
\end{enumerate}

\medskip
Assume that  $X$ is  $(C_m, \alpha_{i+1}, \alpha_{i}), 1 < i+1 < m$. Then we have
$\mathcal C_o(Y)=\mathcal H_{L_0}(\frak l_{-1} \oplus \frak l_{-2})$. By a similar argument, after replacing $\mathcal H_{L_0}(\frak l_{-1})$ by $\mathcal H_{L_0}(\frak l_{-1} \oplus \frak l_{-2})$, we get the desired result. As above, we use the following description of the representation $U_{-1} \oplus \frak l_{-1} \oplus \frak l_{-2}  $
of $L_0$ and $\mathcal H_{L_0}(U_{-1} \oplus \frak l_{-1} \oplus \frak l_{-2} )$.
\begin{enumerate}
\item[$ {\rm (3)_{i+1}}$] $(C_m, \alpha_{i+1}, \alpha_{i})$, $1 < i+1 < m$ where $U=V_{\varpi_1}$; Denote $L_0=A_{i}\times C_{m-i-1}$. Let $V^*$ be the standard representation of $A_{i}$, let $Q^*$ be the standard representation of $C_{m-i-1}$ and $W:= \mathbb C \oplus Q$. Then
\begin{eqnarray*}
\mathfrak l_{-2}&=& \Sym^2V \\
U_{-1} \oplus\mathfrak l_{-1} &=& V \otimes W.
\end{eqnarray*}
%The number $K^{-1}_X.C_{\alpha} -2$   is $2m-(i+1)$. % and is equal to the dimension of
The closure $\mathcal H_{L_0}(U_{-1} \oplus \frak l_{-1} \oplus \frak l_{-2}  )$ of the  $L_0$-orbit at $v\otimes w + v^2 $ is %, where $w$ is a highest weight vector of $W^*$, is
%\begin{eqnarray*} \overline{L_0.(w^2 + w\otimes q +w)} \subset \mathbb P (\Sym^2 W^* +W^*\otimes Q^*+ W^*)=\mathbb P(\mathfrak g_{-2}+\mathfrak g_{-1}),\end{eqnarray*}
\begin{eqnarray*}
\mathcal H_{L_0}(U_{-1} \oplus \frak l_{-1} \oplus \frak l_{-2} ) &=& \mathbb P \{c v \otimes w + v^2 : c \in \mathbb C, v \in V, w \in W \}\\
& \simeq &\mathbb P (\mathcal O_{\mathbb P(V)}(-1)^{2m-2i-1} \oplus \mathcal O_{\mathbb P(V)}(-2)),
\end{eqnarray*}
where $\dim V=i+1$ and $\dim W=2m-2i-1$.
\end{enumerate}
\end{proof}


\section{Smooth horosphrical varieties of type $B_m$ or of type $F_4$} \label{sect:horospherical BF type}

From now on let $X$ be one of the horospherical varieties $(B_m, \alpha_{m-1}, \alpha_{m})$, where $m \geq 3$ or $(F_4, \alpha_2, \alpha_3)$.
Let $\mathfrak g = \bigoplus_{p \in \mathbb Z} \mathfrak g_p $ be the Lie algebra of  $\Aut(X)$ with the  gradation given as in Proposition \ref{tangentsp of X} and  $\mathfrak m:= \bigoplus _{p <0} \mathfrak g_p$ be its negative part.
Set $${\bf U}:= V \oplus (\Sym^2(V) \otimes W)$$ where $V$ and $W$ are vector spaces with $(\dim V, \dim W)=(2, m-1)$ for $(B_m, \alpha_{m-1}, \alpha_{m})$ and $(\dim V, \dim W)=(3, 2)$ for $(F_4, \alpha_2, \alpha_3)$. Then $\bf U$ can be identified with $\frak g_{-1}$. Let ${{\bf S}}  \subset \mathbb P \bf U$ be the variety of minimal rational tangents of $X$ at the base point as in Proposition \ref{vmrt of X}, which is given by
$${\bf S}=\mathbb P \{ v + v^2 \otimes w: v \in V, w \in W\} \subset \mathbb P\bf U.$$

\subsection{Projective geometries of VMRTs}

We will show that  the variety ${{\bf S}}  \subset \mathbb P \bf U$   of minimal rational tangents of $X$ at the base point satisfies the conditions in Proposition \ref{parallel transport of III}.

\begin{lemma} \label{parallel transport of III-X case} %$\,$  %Let $s=\dim \frak g_{-1} -(1+p+r)$.
 Let $s$ be  the codimension of $T^{(2)}\widehat{\bf S}$ in ${\bf U}=\frak g_{-1}$.

  \begin{enumerate}
  \item The third fundamental form of ${\bf S} \subset \mathbb P(\bf U)$ is surjective;
  \item The dimension of $  [\frak g_{-\alpha}, \frak g_{-1}]   $ is   $  s$, where $\alpha$ is the simple root which gives the gradation on $\frak g$.
\end{enumerate}
\end{lemma}

\begin{proof}
 (1) The tangent space $T_{\beta} \widehat{\bf S}$ at $\beta = v + v^2 \otimes w \in \widehat{\bf S} $ is given by
$$T_{\beta}\widehat{{\bf S}} = \{ v' + 2 v \circ v' \otimes w + v^2 \otimes w': v' \in V, w' \in W\}.$$
The second fundamental form $II_{\beta} : \Sym^2 T_{\beta}\widehat{{\bf S}} \rightarrow {\bf U} /T_{\beta}\widehat{{\bf S}}$ is
\begin{eqnarray*}
&& II_{\beta}(v' + 2v \circ v' \otimes w, v'' + 2 v\circ v'' \otimes w) = 2v' \circ v'' \otimes w \\
&& II_{\beta}(v' + 2v \circ v' \otimes w, v^2 \otimes w') = 2 v \circ v' \otimes w' \\
&& II_{\beta}(v^2 \otimes w', v^2 \otimes w'')=0
\end{eqnarray*}
where $v',v'' \in V$ and $w', w'' \in W$. The third fundamental form $III_{\beta} : \Sym^3 T_{\beta}\widehat{{\bf S}} \rightarrow {\bf U} /T_{\beta}^{(2)}\widehat{{\bf S}}$ is zero except
$$III_{\beta}(v' + 2 v\circ v' \otimes w, v'' + 2 v \circ v'' \otimes w, v^2 \otimes w') = 2 v' \circ v'' \otimes w' \mod T_{\beta}^{(2)}\widehat{{\bf S}}, $$
where $v',v'' \in V$ and $w' \in W$.
Thus the third osculating space is the whole space $\bf U$.

(2) The dimension $r$ of the image of $II_{\beta}$ is $m$ in the first case, and is $7$ in the second case. $(\dim{\bf{U}}, \dim\widehat{{\bf S}})=(\dim D_x, p+1)$ is $(2+3(m-1)=3m-1,m+1)$ in the first case, and is $(15, 5)$ in the second case. Thus $s=\dim {\frak g_{-1}} -(1+p+r)$ is $ m-2 $ in the first case, and is $ 3$ in the second case.

On the other hand, the minimum of $\dim [\xi, D_x]$ occurs when $\xi \in \frak g _{-\alpha}$ where $\alpha$ is the simple root which gives the gradation on $\frak g$. Furthermore, the dimension of $[\xi , D_x]=[\xi , \mathfrak l_{-1}]$  is equal to the number of roots $\gamma$ with $\mathfrak g_{-\gamma} \subset  \mathfrak l_{-1}$ such that $\alpha + \gamma$ is again a root. One can compute directly that this number is  $m-2 $ in the first case, and is $3$ in the second case.
\end{proof}

From the exact sequence of $G_0$-modules where $G_0$ is the subgroup of $G=\Aut(X)$ with Lie algebra $\frak g_0$
$$0 \rightarrow V \rightarrow {\bf U} \rightarrow \Sym^2V \otimes W \rightarrow 0, $$
we get the following exact sequences.

\begin{lemma} \label{lem:finer exact sequences} Let $\beta=v + v^2 \otimes w $ be an element of  $ \widehat{\bf S}$, where $v \not=0 \in V$ and $w\not=0 \in W$. Denote by $V_0$ the subspace of $V$ generated by $v$ and by $W_0$ the subspace of $W$ generated by $w$. Then
we have the following exact sequences.
%
\begin{eqnarray*}
0 \rightarrow &\mathbb C \beta& \rightarrow \Sym^2 V_0 \otimes W_0 \rightarrow 0 \\
0 \rightarrow V_0 \rightarrow &T_{\beta}{\widehat{\bf S}}/\mathbb C \beta& \rightarrow \left(V_0 \circ(V/V_0) \otimes W_0\right) \oplus \left( \Sym^2 V_0 \otimes(W/W_0)\right) \rightarrow 0 \\
0 \rightarrow V/V_0 \rightarrow &T^{(2)}_{\beta}\widehat{\bf S}/T_{\beta}\widehat{\bf S}& \rightarrow \left(\Sym^2(V/V_0) \otimes W_0\right) \oplus \left( V_0 \circ (V/V_0) \right) \otimes(W/W_0) \rightarrow 0 \\
0 \rightarrow &{\bf U}/T^{(2)}_{\beta}\widehat{\bf S}& \rightarrow \Sym^2(V/V_0) \otimes (W/W_0) \rightarrow 0
\end{eqnarray*}
\end{lemma}

\begin{proof} From the proof of Lemma \ref{parallel transport of III-X case}, we get the following exact sequences.
\begin{eqnarray*}
0 \rightarrow &\mathbb C \beta& \rightarrow \Sym^2 V_0 \otimes W_0 \rightarrow 0 \\
0 \rightarrow V_0 \rightarrow &T _{\beta}\widehat{\bf S}& \rightarrow V_0 \circ V \otimes W_0 + \Sym^2V_0 \otimes W \rightarrow 0 \\
0 \rightarrow V \rightarrow  &T^{(2)}_{\beta}\widehat{\bf S}& \rightarrow \Sym^2V \otimes W_0 + V_0 \circ V \otimes W \rightarrow 0.
\end{eqnarray*}
Here, we remark that $V_0 \circ V \otimes W_0 + \Sym^2V_0 \otimes W$ and $\Sym^2V \otimes W_0 + V_0 \circ V \otimes W$ are not direct sums:

\begin{eqnarray*}
(V_0 \circ V \otimes W_0)  \cap  (\Sym^2V_0 \otimes W)&=&\Sym^2V_0 \otimes W_0 \\
(\Sym^2V \otimes W_0) \cap ( V_0 \circ V \otimes W)&=& V_0 \circ V \otimes  W_0.
\end{eqnarray*}
%
By taking the quotients we get the desired exact sequences.
  \end{proof}

\begin{remark}\label{pqrst}
Let $p=\dim {\bf S}$, $q=\dim {\bf Q}/  T_{\beta}\widehat{\bf S}$, $r=\dim T^{(2)}_{\beta}\widehat{\bf S}/T_{\beta}\widehat{\bf S}$, $s=\dim {\bf U}/T^{(2)}_{\beta}\widehat{\bf S}$ as above, and $t=q-r-2s$, where ${\bf Q}:= \mathfrak{m}$ as a vector space. Then we have the following table.

\begin{center}
    \begin{tabular}{ r | c c c c c }
 & p & q & r & s & t\\
 \hline
$(B_m, \alpha_{m-1}, \alpha_{m})$ &$m$&$2m-2$&$m$&$m-2$&${(m-2)(m-3)}/{2}$ \\
$(F_{4}, \alpha_{2}, \alpha_{3})$ &$4$&$10$&$7$&$3$&$5$ \\
    \end{tabular}
\end{center}

\end{remark}


\subsection{Parallel transports of VMRTs}
We will show that ${\bf S}$ is not changed under the deformation keeping the second fundamental form and the third fundamental form constant (Proposition \ref{parallel transport of vmrt}).
We adapt arguments in the proof of  Proposition 8.9 of \cite{Hw15}, which proves the same statement as in Proposition \ref{parallel transport of vmrt}  for the case when $X$ is of type $(G_2, \alpha_2,\alpha_1)$. \\

Assume that  $X$ is of type $(F_4, \alpha_2, \alpha_3)$.
Then ${\bf U}$ and ${\bf S}$ are given by   %${\bf S} \subset \mathbb P {\bf U}$ be the variety of minimal rational tangents of $(F_4, \alpha_2, \alpha_3)$.
  ${\bf U} = V \oplus (\Sym^2(V) \otimes W)$ and
\begin{eqnarray*}{\bf S}=\mathbb P \{ c v + v^2 \otimes w: c \in \mathbb C, v \in V, w \in W\},
\end{eqnarray*}
where  $V$ is a vector space of dimension   $3$ and $W$ is a vector space of dimension  $2$.

Consider ${\bf S}$ as a projective bundle $\mathbb P(\mathcal O(-1) \oplus \mathcal O(-2)^2)$ over $\mathbb P(V) $. Let $\psi \colon {\bf S}  \rightarrow \mathbb P(V) $ denote this $\mathbb P^2$-fibration and
let $\xi$ denote the dual tautological line bundle on $\psi:{\bf S} \rightarrow \mathbb P(V)$. Then $H^0({\bf S}, \xi) = H^0(\mathbb P(V), \psi_*\xi) =V^* \oplus (\Sym^2 V ^*\otimes W^*)$. Hence $\xi$ induces the embedding ${\bf S} \subset \mathbb P({\bf U})$.



The $\mathcal O(-1)$-factor defines a subvariety $Z=\mathbb P(\mathcal O(-1)) \subset {\bf S}$, which is isomorphic to $   \mathbb P^2$  %Let $Z \subset {\bf S}$ be the subvariety corresponding to $\mathbb P(\mathcal O(-1)) \simeq \mathbb P^2$.
 and is a section of $\psi$. A choice of an $\mathcal O(-2)$-factor gives a section $B$ of $\psi$ whose linear span $A$ is isomorphic to $\mathbb P^5$ and disjoint from $Z$. We call such a section a \emph{complementary section}. %We have the exact sequence $$ 0 \rightarrow \mathbb P(\mathcal O(-1) \oplus \mathcal O(-2))  \rightarrow  \mathbb P(\mathcal O(-1) \oplus \mathcal O(-2)^2) \rightarrow  B \simeq \mathcal O(-2) \rightarrow 0.$$

Fix an $\mathcal O(-2)$-factor and denote by $B_0$ the corresponding complementary section.
Then the complement ${\bf S} -Z$ of $Z$ is biholomorphic to the total space of the vector bundle $\mathcal O(1)^2$ on $\mathbb P(V)$ whose zero section corresponds to $B_0$. A complementary section corresponds to a section of $\mathcal O(1)^2$. Thus,  for a given   triple $(s_1^1, s_1^2, s_1^3)$ in ${\bf S} -Z$  such that $\psi(s_1^1), \psi(s_1^2), \psi(s_1^3)$ are distinct,  there is a unique complementary section  $B_1$   with  $s_1^1, s_1^2, s_1^3 \in B_1$.

\begin{proposition} \label{characterization of vmrt F}  When $X$ is of type $(F_4, \alpha_2, \alpha_3)$,
 ${\bf S} \subset \mathbb P^{14}$ is obtained from the following four data:
   \begin{enumerate}
\item[(i)] a plane $Z \subset \mathbb P^{14}$,
\item[(ii)] two disjoint linear spaces $A_1 \cong \mathbb P^5$ and $A_2 \cong \mathbb P^5$
such that the union $Z \cup A_1 \cup A_2$ span $\mathbb P^{14}$, %whose span is disjoint from $Z$, so that  $\mathbb P(\hat Z \oplus \hat A_1 \oplus \hat A_2) \cong \mathbb P^{14}$,
\item[(iii)] two subvarieties $B_1 \subset A_1$ and $B_2 \subset A_2$, each of which is  isomorphic to the Veronese surface $\nu(\mathbb P^2) \subset \mathbb P^5$,
\item[(iv)] two birational maps $\epsilon_1 \colon Z \rightarrow B_1$ and $\epsilon_2 \colon Z \rightarrow B_2$.
   \end{enumerate}
\end{proposition}


\begin{proof}
The locus of planes generated by $z \in Z$ and $\varepsilon_1(z), \varepsilon_2(z)$ is a subvariety of $\mathbb P^{14}$ projectively equivalent to ${\bf S} \subset \mathbb P^{14}$.
\end{proof}


Similarly we have the following.

\begin{proposition} \label{characterization of vmrt B}  When $X$ is of type $(B_m, \alpha_{m-1}, \alpha_m)$,
 ${\bf S} \subset \mathbb P^{ 3m-2}$ is obtained from the following four data:
   \begin{enumerate}
\item[(i)] a line $Z \subset \mathbb P^{3m-2}$,
\item[(ii)] $m-1$ linear spaces $A_i \cong \mathbb P^3$, where $1 \leq i \leq m-1$, such that the union $Z \cup \left(\cup_{i=1}^{m-1} A_i\right) $ spans $ \mathbb P^{3m-2}$,
\item[(iii)] $m-1$ subvarieties $B_i \subset A_i$, where $1 \leq i \leq m-1$, each of   which is  isomorphic to the conic $\nu(\mathbb P^1) \subset \mathbb P^2$,
\item[(iv)] $m-1$ birational maps $\epsilon_i \colon Z \rightarrow B_i$, where $1 \leq i \leq m-1$.
   \end{enumerate}
\end{proposition}


\begin{proposition} \label{parallel transport of vmrt}
Let ${\bf S} \subset \mathbb P {\bf U} $ be the variety of minimal rational tangents of
$(B_m, \alpha_{m-1}, \alpha_{m})$, $m \geq 3$ or $(F_4, \alpha_2, \alpha_3)$ at the base point.


  Let $\pi: \mathbb P \mathcal U\rightarrow \mathbb P^1$ be the projectivization of a holomorphic vector bundle $\mathcal U$ over  $  \mathbb P^1$  and let $\mathcal C \subset \mathbb P \mathcal U$ be an irreducible subvariety. Denote by $\varpi$ the restriction of $\pi$ to $\mathcal C$. Assume that
   \begin{enumerate}
   \item $\mathcal C_t:=\varpi^{-1}(t) \subset \mathbb P  \mathcal U_t:=\pi^{-1}(t)$ is projectively equivalent to ${\bf S} \subset \mathbb P {\bf U}$ for all $t \in \mathbb P^1 -\{t_1, \dots, t_k\}$; % but (possibly) finitely many points of $\mathbb P^1$;
  \item
      for  a general  section $\sigma \subset \mathcal C$ of $\varpi$, the relative second fundamental forms and the relative third fundamental forms of $\mathcal C$ along $\sigma$ are constants.
     \end{enumerate}

Then for any $t \in \mathbb P^1$, $\mathcal C_t \subset \mathbb P(\mathcal U_t)$ is projectively equivalent  to ${\bf S} \subset \mathbb P({\bf U})$.
\end{proposition}


\begin{proof} Assume that $\mathcal C_t \subset \mathbb P \mathcal  U_t$ is projectively equivalent to ${\bf S} \subset \mathbb P {\bf U}$ for all $t  $ in the unit disc $\Delta \subset \mathbb P^1$ except  for $t=0$. Assume further  that for  a general  section $\sigma \subset \mathcal C$ of $\varpi$, the relative second fundamental forms and the relative third fundamental forms of $\mathcal C$ along $\sigma$ are constants. Then there is an open submanifold $\mathcal C^0\subset \mathcal C$ such that $\mathcal C_t^0 \subset \mathcal C_t $ corresponds to an open subset of ${\bf S} -Z \subset {\bf S}$ for $t \not=0  \in \Delta $.   We claim that $\mathcal C_0 \subset \mathbb P \mathcal U_0$ is also projectively equivalent to ${\bf S} \subset \mathbb P {\bf U}$.
%
It suffice to show that we   get   the four data in Proposition \ref{characterization of vmrt F} and Proposition \ref{characterization of vmrt B} at $t=0$ from $\mathcal C_0=\varpi^{-1}(0) \subset \mathbb P \mathcal  U_0$.
\\

Assume that $X$ is of type $(F_4, \alpha_2, \alpha_3)$.
Let $\beta_t \in \mathcal C_t$  be a section of $\varpi$ corresponding to a general point $\beta \in {\bf S}$ for $t\neq 0$ and let $\beta_0$ be the limit. According to the computation of the second fundamental form in the proof of  Lemma  \ref{parallel transport of III-X case} (1), we see that $Baselocus(II_{\beta_t})$ gives us a foliation whose leaves are isomorphic to $\mathbb P^2$-fiber of the projective bundle ${\bf S}$ over $\mathbb P^2$. By the assumption (2), there is also a foliation on the central fiber given by $Baselocus(II_{\beta_0})$, which is exactly $\mathbb P^2$ in $\mathcal C_0=\varpi^{-1}(0) \subset \mathbb P \mathcal  U_0$. We call these subvarieties $\mathbb P^2$ of $\mathcal C_t$ for $t \in \Delta$, corresponding to the $\mathbb P^2$-fiber of  ${\bf S}$, \emph{the planes of the rulings}. \\

(i) Let $\psi_t \colon \mathcal C_t  \rightarrow \mathbb P^2$ be the coresponding fibration for $t \neq 0$ of the projective bundle ${\bf S}=\mathbb P(\mathcal O(-1) \oplus \mathcal O(-2)^2) \rightarrow \mathbb P^2$. Then $Z \subset {\bf S}$ gives familly of distingushied planes $Z_t \simeq \mathbb P(\mathcal O(-1))) \subset \mathcal C_t \subset \mathbb P  \mathcal U_t$ for $t\neq 0$, which is a section of $\psi_t$. Then the limit $Z_0$ is also a $\mathbb P^2$. Hence, there are $\mathbb P^2$ subbundle  $\mathcal Z \subset \mathcal C$ of $\pi: \mathbb P \mathcal U\rightarrow \mathbb P^1$. \\

(ii) Pick 3 distinct points $s_1^1, s_1^2, s_1^3 \in \mathcal C_0 \backslash Z_0$ that lie  in 3 distinct planes of  the rulings and choose local sections $\sigma^1_1, \sigma^2_1, \sigma^2_1$ of $\varpi^0:\mathcal C^0 -\mathcal Z \rightarrow \Delta$ such that $\sigma^i_1(0)=s_1^i$ for $i=1, 2, 3$. Then there exists a unique complementary section $B_{1,t} $
with $\sigma^1_1(t), \sigma^2_1(t), \sigma^3_1(t) \in B_{1,t}$  for any $t \not=0$.

Let  $A_{1,t}$ be the linear span of $B_{1,t}$.
Then,   $B_{1,t} \subset A_{1,t}$ is isomorphic to the Veronese surface $\nu(\mathbb P^2) \subset \mathbb P^5$. The limit $A_{1,0}$ of $A_{1,t}\simeq \mathbb P^5 \subset \mathbb P \mathcal U_t$ is  also a projective space  $  \mathbb P^5 \subset \mathbb P \mathcal U_0$, which contains $s_1, s_2, s_3  $.

Choose 3 distinct points  $s^1_2, s^2_2, s^3_2 $ in the complement $\mathcal C_0 - (Z_0 \cup A_{1,0})$ such that  each line $span(s^i_1, s^i_2)$ for $i=1,2,3$ is in the same  plane of the rulings. Choose local sections $\sigma^i_2$ of $\varpi$ such that $\sigma^i_2(t) \in \mathcal C_t  - (Z_t \cup  A_{1,t})$ and $\sigma^i_2(0)=s_2^i$ for $i=1, 2, 3$. Then there exist a  unique complementary section $B_{2,t} $

 such that $\sigma^1_2(t), \sigma^2_2(t), \sigma^3_2(t) \in B_{2,t}$.

 Let  $A_{2,t}$ be the linear span of $B_{2,t}$. % isomorphic to $\mathbb P^5$.
 The two complementary sections $B_{1,t} \subset A_{1,t}$ and $B_{2,t} \subset A_{2,t}$ are isomorphic to the Veronese surface $\nu(\mathbb P^2) \subset \mathbb P^5$ for $t \not=0$.

 The limit $A_{j,0}$ of $A_{j,t}\simeq \mathbb P^5 \subset \mathbb P \mathcal U_t$  is also projective space $ \mathbb P^5 \subset \mathbb P \mathcal U_0$, $j=1, 2$ such that $s^1_1, s^2_1, s^3_1 \in A_{1,0}$ and $s^1_2, s^2_2, s^3_2 \in A_{2,0}$.
 Then the union $Z_0 \cup A_{1,0} \cup A_{2,0}$ spans $\mathbb P \mathcal U_0$. For, otherwise, there is hyperplane such that $Z_0 \cup A_{1,0}  \cup A_{2,0} \subset \mathbb P^{13} \subset \mathbb P \mathcal U_0$ which contradicts to the constancy of second and third fundamental form. \\

(iii) Consider the limit $B_{j.0}$ of $B_{j,t}$, which is a surface in $A_{j,0}=\mathbb P^5$ of degree less then $4$. The planes of rulings on $\mathcal C_0$ intersecting $A_{j,0}$, give  a analytic surface $B'_{j,0} \subset B_{j,0}$. Let $A'_{j,0}$ be the linear span of $B'_{j,0}$. If $B'_{j,0}$ lies in a hyperplace of $A_{j,0}$, then $\mathbb P (Z_0 \oplus A'_{1,0}\oplus A'_{2,0})$ is contained in a hyperplane of $\mathbb P U_0$ which contradiction to the fundamental forms. Hence, $B'_{j,0}$ must be non-degenerate in $A_{j,0}$ and this means $B_{j,0}$ is irreducible non-degenerate surface. Since the Veronese surface is minimal degree surface, a Veronese surface for $j=1,2$. \\

(iv) With a choice of $Z_t$, $B_{j,t}$ and $A_{1,t}$ for $j=1, 2$ as above in (i) - (ii), we have birational morphisms $\varepsilon_{j,t}\colon Z_0 \rightarrow B_{j,t}$ for $t \neq 0$ such that the planes of rulings on $\mathcal C_t$ are the plane spanned by $z_t \in Z_t$, $\varepsilon_{1,t}(z_t)  \in B_{1,t}$ and $\varepsilon_{2,t}(z_t)  \in B_{2,t}$.
Let $\tilde{\mathcal P}$ be the blow up of the bundle $\pi: \mathbb P \mathcal U\rightarrow \Delta$ along the submanifold $\mathcal Z$ and let ${\mathcal E}$ be the exceptional divisor,  which is biholomorphic to $\mathcal Z \times _{\Delta}\mathbb P (\widehat {\mathcal A_1} \oplus \widehat{\mathcal A_2})$. Let $E_t$ be the exceptional divisor of the blow up of $\mathbb P \mathcal U_t$, which is $Z_t \times \mathbb P (\widehat{A_{1,t}} \oplus \widehat{A_{2,t}})$. % as a bundle over $\mathbb P^1$.
 For $t \neq 0$, the proper transform of the plane joining $z \in Z_t$, $\varepsilon_{1,t}(z)  \in B_{1,t}$ and $\varepsilon_{2,t}(z)  \in B_{2,t}$ intersects $E_t$ at a point $(z,\epsilon_t(z)) \in  E_t$, where $\epsilon_t(z):=\varepsilon_{1,t}(z)\oplus\varepsilon_{2,t}(z)$. The surface $\Gamma_t:=\{(z, \varepsilon_t (z)| z \in Z_t\}$ corresponds to the graph of $\epsilon_t$.
Let $\Gamma_0 \subset E_0$ be the limit of $\Gamma_t$.
%
Then there exists an irreducible component $\Gamma'_0$ of $\Gamma_0$ which is the intersection of $E_0$ with the proper transform of the planes in the planes of ruling at $t=0$. By definition, $\Gamma'_0$ dominate $B_{j,0}$ for $j=1, 2$. $\Gamma'_0$   dominates $Z_0$, too. For, otherwise, the intersection with the planes of ruling of $Z_0$ is codimension $\geq 1$, which contradicts to the constancy of the second fundamental forms and the third fundamental forms.

On $\mathcal E$, let $\mathcal L^Z$ be the line bundle which is the pull back of the relative $\mathcal O(1)$ bundle on $\mathcal Z$ and  let $\mathcal L^{A_j}$ be the line bundle which is the pull back of the relative $\mathcal O(1)$ bundle on $\mathcal A_{j}$. Then $\mathcal L^Z \otimes \mathcal L^{A_1} \otimes \mathcal L^{A_2}$ be a relative ample line bundle on $\mathcal E$ which has degree 9 with respect to $\Gamma_t$ for $t\neq 0$. Since $\Gamma'_0$ dominate $Z_0$ and $B_{j,0}$ for $j=1, 2$, we have $\Gamma'_0\cdot \mathcal L^{Z}\geq 1$ and $\Gamma'_0\cdot \mathcal L^{A_{j}} \geq 4$. Hence, $\Gamma'_0=\Gamma_0$ and $\Gamma_0$ determines birational map $\varepsilon_{j,0} \colon Z_0 \rightarrow B_{j,0}$ for $j=1, 2$. \\

In case $(B_m, \alpha_{m-1}, \alpha_{m})$, $m>2$,
\begin{eqnarray*}{\bf S}=\mathbb P \{ c v + v^2 \otimes w: c \in \mathbb C, v \in V, w \in W\} \simeq \mathbb P(\mathcal O(-1) \oplus \mathcal O(-2)^{m-1})\end{eqnarray*}
where $\dim V =2$ and $\dim W= m-1$. Since $ \nu_2 (\mathbb P^1) \subset \mathbb P \Sym^2 V \simeq  \mathbb P^2$ is rational normal curve with minimal degree, we conclude the result in a similar method.
\end{proof}



\subsection{Fundamental graded Lie algebras determined by VMRTs}    \label{sect: geometric structures modeled after X}

We will show that ${\bf S}$ and $\frak m$ satisfies the conditions in
Proposition \ref{S-str and G_0-str}, so that there is a one-to-one correspondence between $G_0$-structures and ${\bf S}$-structures on filtered manifolds of type $\frak m$.

\begin{proposition} \label{fundamental graded Lie algebra of vmrt}  $\,$
 \begin{enumerate}
\item The graded Lie algebra  $\mathfrak m $ is fundamental and determined by ${\bf S}  \subset \mathbb P \mathfrak g_{-1}$.
\item Let $\frak g(\widehat{\bf S})$ be the Lie algebra of the linear automorphism group   $G(\widehat{\bf S})$.
Then the induced homomorphism $\frak g(\widehat{\bf S}) \rightarrow \frak g_0(\frak m)$ is injective and its image is $\frak g_0$.
\end{enumerate}
 \end{proposition}

In the remaining part of this subsection we will prove Proposition \ref{fundamental graded Lie algebra of vmrt}. % and Proposition \ref{extension to rationally connected open subset}.
Recall that ${\bf U}$ is given by  $${\bf U} = V \oplus (\Sym^2(V) \otimes W)$$ where $V$ and $W$ are vector spaces with $(\dim V, \dim W)=(2, m-1)$ for $(B_m, \alpha_{m-1}, \alpha_{m})$ and $(\dim V, \dim W)=(3, 2)$ for $(F_4, \alpha_2, \alpha_3)$ and  that    ${{\bf S}}$ is given by
$${\bf S}=\mathbb P \{ v + v^2 \otimes w: v \in V, w \in W\}  \simeq \mathbb P (\mathcal O(-1) \oplus \mathcal O(-2)^{k}) \subset \mathbb P\bf U,$$ where $k=\dim W$.
Let ${\bf G}:=SL(V) \times SL(W)$ be a group acting on $\bf S$. Then, ${\bf S}$ is a smooth horospherical ${\bf G}$-variety of rank one. Moreover,  it has two closed ${\bf G}$-orbits, ${\bf Z}  :=\mathbb P(V)$ corresponding to $\mathbb P(\mathcal O(-1))$ and
\begin{eqnarray*}  {\bf Y}   := \mathbb P \{\lambda^2\otimes\mu \in \Sym^2V\otimes W | \lambda \in V ,
\mu \in W \}\end{eqnarray*}
corresponding to a choice of  $\mathbb P(\mathcal O(-2)^{k})$, has one open ${\bf G}$-orbit the complement of ${\bf Y} \cup {\bf Z}$ in ${\bf S}$.

Denote that ${\bf Y} $ is the variety of minimal rational tangents of $Y=L/P_Y$ that is a rational homogeneous associated with a long simple root of type $(B_m, \alpha_{m-1})$ (respectively, of type $(F_4, \alpha_2)$).

\begin{proposition} [Proposition 1 and Proposition 7 of \cite{HM02}, Proposition 4.1 of \cite{HH}] \label{Vmrt ss-LieAl Str}  Let $Y$ be a rational homogeneous space of type $(\mathfrak l , \alpha)$ where  $\alpha$ is  a long simple root. Let $\mathfrak n=\bigoplus _{p <0} \frak l_p$ be the negative part of the graded Lie algebra  $\frak l=\bigoplus_{p\in \mathbb Z}\frak l_p$ with gradation associated to $\alpha$. Let $D_o \subset T_oY$ be the linear span of the homogeneous cone of the variety $\mathcal C_o$ of minimal rational tangents  at a base point $o \in Y$. Then $\frak n$ is the fundamental graded Lie algebra determined by $ \mathcal C_o \subset \mathbb PD_o$.
\end{proposition}


Since $\wedge^2 (\Sym^2 V \otimes W) = \wedge^2 (\Sym^2 V) \otimes \Sym^2 W \oplus \Sym^2 (\Sym^2 V) \otimes \wedge^2 W$, if we decompose $\Sym^2(\Sym^2 V) = \Sym^4V \oplus (\Sym^4V)^{\perp}$, then the Lie bracket $[\,,\,]:\wedge^2 \frak l_{-1} \rightarrow \frak l_{-2}$ is given by the projection map
\begin{eqnarray*}
  \nu:  \wedge^2 (\Sym^2 V \otimes W) & \to &(\Sym^4V)^{\perp} \otimes \wedge^2 W.
\end{eqnarray*}
If $X$ is of type $(B_m, \alpha_{m-1}, \alpha_m)$, then $(\Sym^4V)^{\perp} \otimes \wedge^2 W =  \mathbb C  \otimes \wedge^2 W =  \wedge^2 W$. If $X$ is of type $(F_4, \alpha_2, \alpha_3)$, then $(\Sym^4V)^{\perp} \otimes \wedge^2 W =   \Sym^2(\wedge^2V)  \otimes \mathbb C=\Sym^2(\wedge^2V)$.
Thus $\wedge^2 (\Sym^2 V \otimes W)$ is the direct sum of $\Ker \nu$ and an irreducible representation of $SL(V) \times SL(W)$.

Note that $\frak g_{-1} ={\bf U}$ and $ \frak g_{-2} =  \frak l_{-2}$. The  Lie bracket $[\,,\,]:\wedge^2 \frak g_{-1} \rightarrow \frak g_{-2}$ defines
$$\omega:\wedge^2 (V\oplus \Sym^2 V \otimes W)  \rightarrow \frak l_{-2},$$  where $\omega|_{\wedge^2 (\Sym^2 V \otimes W)}=\nu$ and $\omega(V,V)= \omega(V, \Sym^2V \otimes W)=0$. Then,  $\wedge^2 {\bf U}=\wedge^2 (V\oplus \Sym^2 V \otimes W)$ is decomposed as
$$\wedge^2 (V\oplus \Sym^2 V \otimes W) = \wedge^2 V\oplus (V \wedge (\Sym^2 V\otimes W))\oplus \wedge^2(\Sym^2 V\otimes W) $$
and we have   $\Ker \omega = \wedge^2 V \oplus (V \wedge (\Sym^2 V\otimes W))\oplus \Ker \nu$ %under the decomposition



\begin{lemma}\label{tangent space of the cone V-2}
The kernel $\Ker \omega $ is spanned by $\wedge^2 P \subset \wedge^2({\bf U})$ where  $P$ is 2-dimensional subspace of ${\bf U}$ tangent to $\widehat{ \bf S }$.
\end{lemma}

\begin{proof} Let $\Xi$ be the subspace of $\wedge^2 (V\oplus \Sym^2 V \otimes W)$ spanned by
   \begin{eqnarray*} \{ \wedge^2 P \subset \wedge^2(V\oplus \Sym^2V \otimes W) | P \mbox{ is tangent to } \widehat{ \bf S }, \dim P=2 \}. \end{eqnarray*}

Since $\widehat{ \bf S } \subset V\oplus \Sym^2V \otimes W $ is a $SL(V)\times SL(W)$-invariant subvariety,

$\Xi$  is also $SL(V)\times SL(W)$-invariant subspace of $\wedge^2 (V\oplus \Sym^2 V \otimes W)$. Furthermore, $\wedge^2 (V\oplus \Sym^2 V \otimes W)$ is the direct sum of $\Ker \omega$ and an irreducible representation of $SL(V) \times SL(W)$.
We will show that $\Ker \omega  $ is contained in $ \Xi$. If so, since $\Xi$ is an $SL(V) \times SL(W)$-invariant subspace and is proper, we get the desired equality.


As the subspace $V$ is contained in $\widehat{\bf S}$, the component $\wedge^2 V$ is contained in $\Xi$.
By Proposition \ref{Vmrt ss-LieAl Str}, the kernel of $\nu$ is spanned by $\wedge^2 P$, where $P$ is tangent to $\widehat{\bf Y}$, and thus is contained in $\Xi$.
We claim that  $V \wedge (\Sym^2 V \otimes W)$ contained in $\Xi$, which complete the proof.

Let $\beta_t = v_t + v_t^2 \otimes w_t$ be a curve in $\widehat{\bf S}$, where $v_t \in V$ and $w_t \in W$.  By definition of $\Xi$,
 \begin{eqnarray*}
 \beta_0 \wedge \beta'_0 = (v_0 + v_0 ^2 \otimes w_0) \wedge (v_0' +  v'_0 \otimes v_0 \otimes w_0 + v_0 \otimes v'_0 \otimes w_0 + v_0^2 \otimes w_0')
 \end{eqnarray*}
is contained in $\Xi$. Since $\alpha_t:=v_t^2 \otimes w_t$ is a curve in $\widehat{\bf Y} \subset \widehat{\bf S}$,
 $$\alpha_0 \wedge \alpha'_0=(v_0 ^2 \otimes w_0) \wedge (v'_0 \otimes v_0 \otimes w_0 + v_0 \otimes v'_0 \otimes w_0+ v_0^2 \otimes w_0') $$ is contained in $\Xi$. Hence,
$\beta_0 \wedge \beta'_0 -\alpha_0 \wedge \alpha'_0$ which is equal to
 \begin{eqnarray*}
   v_0  \wedge ( v'_0 \otimes v_0 \otimes w_0 + v_0 \otimes v'_0 \otimes w_0 + v_0^2 \otimes w_0')+ (v_0 ^2 \otimes w_0) \wedge  v_0'
 \end{eqnarray*}
 is contained in $\Xi$. Note that
 \begin{eqnarray*}
&&v_0 \wedge (v_0' \otimes v_0 + v_0 \otimes v_0') \\
&& = v_0 \otimes v_0' \otimes v_0 + v_0 \otimes v_0 \otimes v_0' - v_0' \otimes v_0 \otimes v_0 - v_0 \otimes v_0' \otimes v_0 \\
 &&=  v_0^2 \wedge v_0'
 \end{eqnarray*}
 and
 \begin{eqnarray*}
 v_0 \wedge v_0^2 = v_0 \otimes (v_0 \otimes v_0) -(v_0 \otimes v_0) \otimes v_0 = 0
 \end{eqnarray*}
 Therefore, $(v_0' \wedge v_0^2)\otimes w_0$ is contained in $\Xi$.  Here, we consider $V \wedge (\Sym^2 V \otimes W)$ as $(V \wedge \Sym^2 V )\otimes W$.
 This is true for arbitrary $v_0, v_0' \in V$ and $w_0 \in W$, it follows that $V \wedge (\Sym ^2 V \otimes W)$ is contained in $\Xi$.
\end{proof}


\begin{lemma} \label{automorphism group of S} $\Aut^0({\bf S}) = (SL(V) \times SL(W)) \times \mathbb C^* \vartriangleright V^* \otimes W^*$
\end{lemma}
\begin{proof}% [Proof of Lemma \ref{automorphism group of S}]
Since ${\bf S}$ is a smooth horospherical ${\bf G}$-variety of rank one with two closed ${\bf G}$-orbits ${\bf Z}$ and ${\bf Y}$, we will apply the same arguments as in the proof of Lemma 1.1. of \cite{Pa}.
\begin{center}
    \begin{tabular}{ c | c c c c l }
type &$\dim {\bf S}$ & $\dim {\bf Y}$ & $\dim {\bf Z}$ & $N_{{\bf Y}|{\bf S}}$ & $N_{{\bf Z}|{\bf S}}$ \\
 \hline
$(B_m, \alpha_{m-1}, \alpha_{m})$ &$m$&$m-1$&$1$& $\mathcal O(1)$&$\mathcal O(-1)^{m-1}$ \\
$(F_{4}, \alpha_{2}, \alpha_{3})$ &$4$&$3$&$1$&$\mathcal O(1)$&$\mathcal O(-1)^{2}$ \\
    \end{tabular}
\end{center}
Thus $H^0({\bf Y}, N_{{\bf Y}|{\bf S}}) =V^* \otimes W^*$ and $H^0({\bf Z}, N_{{\bf Z}|{\bf S}})=0$.

By the same arguments as in the proof of Lemma 1.1. of \cite{Pa}, the closed $\bf G$-orbit ${\bf Z} $ is stable under the action of $\Aut^0({\bf S})$ and we have
\begin{eqnarray*}
\Aut^0({\bf S}) &=& ({\bf G} \times \mathbb C^*)  \vartriangleright H^0({\bf Y}, N_{{\bf S}_Y/{\bf S}}) \\
&=& (SL(V) \times SL(W) \times \mathbb C^*)  \vartriangleright V^* \otimes W^*.
\end{eqnarray*}
\end{proof}

\begin{proof} [Proof of Proposition \ref{fundamental graded Lie algebra of vmrt} (1)]
Proposition \ref{Vmrt ss-LieAl Str} says that $\frak l_-=\oplus_{k<0}\frak l_k$ is the fundamental graded Lie algebra determined by ${\bf Y}   \subset \mathbb P(\frak l_{-1})$.
Since $\frak m = U_{-1} \oplus \frak l_-$ and $[U_{-1}, \frak m]=0$,  By Proposition \ref{Vmrt ss-LieAl Str} and Lemma \ref{tangent space of the cone V-2},  the graded Lie algebra $\frak m$ is  fundamental and determined by ${\bf S} \subset \mathbb P\frak g_{-1}$.
\end{proof}

\begin{proof} [Proof of Proposition \ref{fundamental graded Lie algebra of vmrt} (2)]
 We recall that $\mathfrak g_0 =(\mathfrak l_0 \oplus \mathbb C   )\vartriangleright U_{0} \subset \mathfrak g_0(\mathfrak m)$ from Proposition \ref{tangentsp of X}. The center of $\frak l_0$ is of dimension one, the semisimple part of $\frak l_0$ is $\frak{sl}(V) + \frak {sl}(W)$ and the vector space $U_{0}$ is $V^* \otimes W^*$. We compare $\frak g_0$ with the Lie algebra of the neutral component $\Aut^0(\widehat{{\bf S}})$ of the automorphism group of the cone $\widehat{\bf S} \subset \frak g_{-1}$ over $\bf S$. %From Lemma \ref{automorphism group of S}, $\Aut(\widehat{{\bf S}})$ is equal to the linear automorphism group $G(\widehat{\bf S})$  and the induced map $\frak g(\widehat{\bf S}) \rightarrow \frak g_0(\frak m)$ is injective whose image in $\mathfrak g_0(\frak m)$ agrees with $\mathfrak g_0   \subset \mathfrak g_0(\mathfrak m)$.
 Then the rest of proof follows from Lemma \ref{automorphism group of S}, that is, $\Aut^0(\widehat{{\bf S}})$ is equal to the linear automorphism group $G(\widehat{\bf S})$    and the induced map $\frak g(\widehat{\bf S}) \rightarrow \frak g_0(\frak m)$ is injective whose image   in $\mathfrak g_0(\frak m)$ agrees with $\mathfrak g_0   \subset \mathfrak g_0(\mathfrak m)$.
%
\end{proof}

\section{$H^2$-cohomologies} \label{sect:H2 cohomology}


In Section \ref{sect:computation of H2 cohomology} -- \ref{sect: proof of proposition}, we assume that $X$ is either of type $(B_m, \alpha_{m-1}, \alpha_m)$ where $m \geq 3$ or of type $(F_4, \alpha_2, \alpha_3)$ and we keep the same notation as in Section \ref{sect:horospherical BF type}. We write $\mathbb C$ as $\frak z$ to emphasize that $\mathbb C$ commutes with $\mathfrak l$.

\subsection{Reductions}   \label{sect:computation of H2 cohomology}
%
A main ingredient is that the computation of $H^2(\frak m, \frak g)$ can be reduced to that of  $H^2(\frak l_-, \frak g)$ and $H^1(\frak l_-, \frak g)$, which can be computed  by applying Kostant's theory.

\begin{proposition} \label{prop: computation of H2 cohomology} Let $k \geq 1$.
%
 Let $\phi \in \Hom(\wedge^2 \frak m, \frak g)_k$ be such that $\partial \phi=0$.
 Then there is $\eta \in  \Hom(\frak m, \frak g)_k$ and  $\zeta \in \Hom(\wedge^2 \frak g_{-1}, U )_k$  such that
 $$\phi = \partial \eta + \zeta   $$
 and $\zeta|_{\wedge^2 \frak l_-} \in H^2(\frak l_-, U  )_k$ and $\zeta(Y^{-1}, \,\cdot\,)|_{\frak l_{-}} \in H^1(\frak l_-, U)_{k-1}+H^1(\frak l_-, \frak l)_{k-1}$ for any $Y^{-1} \in U_-$ and $\zeta|_{\wedge^2U_-} \in H^2(U_-, U)_1$.
 Furthermore, the choice of $\zeta$ in the  expression $\phi = \partial \eta + \zeta   $ is unique.
\end{proposition}

 We postpone the proof of  Proposition \ref{prop: computation of H2 cohomology} until Section \ref{sect: proof of proposition} and proceed with the computation of the  cohomology $H^2(\frak m, \frak g)$.

\begin{lemma} [Computation of $H^1(\frak l_-, \frak g)$] \label{lem:H1 vanishing semisimple} Let $k \geq 1$.

\begin{enumerate}
\item[(i)] $H^1(\frak l_-, \frak l)_{k-1}$ vanishes except for $k=1$ of type $(B_3, \alpha_{2}, \alpha_3 )$ and $$H^1(\frak l_-, \frak l)_{0} \subset \Hom(\frak l_{-1}, \frak l_{-1}).$$
\item[(ii)] $H^1(\frak l_-, \frak z)_{k-1}$ vanishes except for $k =2$ and $H^1(\frak l_-,\frak z)_1 \subset \Hom(\frak l_{-1}, \frak z)$.
\item[(iii)] $H^1(\frak l_-, U)_{k-1}$ vanishes except for $k =1$ and $$H^1(\frak l_-, U)_0\subset \Hom( \frak l_{-1}, U_{-1}).$$
\end{enumerate}

\end{lemma}

\begin{lemma} [Computation of $H^2(\frak l_-, \frak g)$] \label{lem:H2 vanishing semisimple}
 Let $k \geq 1$.

\begin{enumerate}

\item[(i)] $H^2(\frak l_-, \frak l)_k $ vanishes.

\item[(ii)] $H^2(\frak l_-, \frak z)_k  $ vanishes except for $k=2$, and
$$H^2(\frak l_-, \frak z)_2 \subset \wedge^2 \frak l_{-1}^* \otimes \frak z \subset (\wedge^2 \frak l_{-1} \otimes U_{-1})^*\otimes U_{-1} $$

\item[(iii)]
\begin{enumerate}
\item[(1)] If $(\frak m,\frak g_0)$ is of type $(B_m, \alpha_{m-1}, \alpha_m )$, then  $H^{ 2}(\frak l_-, U)_k$  vanishes except for $k=1,2$,  and we have
    \begin{eqnarray*}
    H^2(\frak l_-,  U)_1 &\subset& \wedge^2 \frak l_{-1}^* \otimes U_{-1} \\
    H^2(\frak l_-, U)_2 & \subset& \wedge^2 \frak l_{-1}^* \otimes U_0 \subset (\wedge ^2 \frak  l_{-1}  \otimes \mathfrak l_{-1})^* \otimes U_{-1} \\
    \end{eqnarray*}
\item[(2)] If $(\frak m,\frak g_0)$ is of type $(F_4, \alpha_2, \alpha_3 )$, then $H^{ 2}(\frak l_-, U)_k$   vanishes except for $k=1$, and we have
    \begin{eqnarray*}
    H^2(\frak l_-, U)_1 &\subset& \wedge^2 \frak l_{-1}^* \otimes U_{-1}
    \end{eqnarray*}
\end{enumerate}
\end{enumerate}
\end{lemma}
\begin{proof} Use Kostant's theory (\cite{Ko61}).

\end{proof}

We remark that, in the proof of Proposition \ref{prop: computation of H2 cohomology}, we use the property that both $H^1(\frak l_-, \frak l)_{k-1}$ and $H^2(\frak l_-, \frak l)_k$ vanish for any $k \geq 1$ (Lemma \ref{lem:H1 vanishing semisimple} (i) and Lemma \ref{lem:H2 vanishing semisimple} (i)).

\begin{proposition} [Computation of $H^2(\frak m, \frak g)$]  \label{prop:H2 vanishing}
 Let $k \geq 1$.  %$\,$
\begin{enumerate}
\item If $(\frak m,\frak g_0)$ is of type $(B_3, \alpha_{2}, \alpha_3 )$,  then $H^{ 2}(\frak m, \frak g)_k$ vanishes except for $k=1,2$, and
    \begin{eqnarray*}
    H^2(\frak m, \frak g)_1 &\subset& \wedge^2 \frak g_{-1}^* \otimes \frak g_{-1} \\
    H^2(\frak m, \frak g)_2 &\subset& \wedge^2 \frak g_{-1}^* \otimes U_0  \subset (\wedge ^2 \frak g_{-1}  \otimes \mathfrak g_{-1})^* \otimes U_{-1}.
    \end{eqnarray*}
\item If $(\frak m,\frak g_0)$ is of type $(B_m, \alpha_{m-1}, \alpha_m )$, where $m >3$,  then $H^{ 2}(\frak m, \frak g)_k$ vanishes except for $k=1,2$, and
    \begin{eqnarray*}
    H^2(\frak m, \frak g)_1 &\subset& \wedge^2 \frak g_{-1}^* \otimes U_{-1} \\
    H^2(\frak m, \frak g)_2 &\subset& \wedge^2 \frak g_{-1}^* \otimes U_0  \subset (\wedge ^2 \frak g_{-1}  \otimes \mathfrak g_{-1})^* \otimes U_{-1}.
    \end{eqnarray*}
\item If $(\frak m,\frak g_0)$ is of type $(F_4, \alpha_2, \alpha_3 )$, then $H^{ 2}(\frak m, \frak g)_k$ vanishes except for $k=1 $, and
  %  $$H^2(\frak m, \frak g)_1 \subset \wedge^2 \frak g_{-1}^* \otimes U_{-1} .$$
    \begin{eqnarray*}
    H^2(\frak m, \frak g)_1 &\subset& \wedge^2 \frak g_{-1}^* \otimes U_{-1}% \\
   % H^2(\frak m, \frak g)_2 &\subset& \wedge^2 \frak g_{-1}^* \otimes \frak z \subset (\wedge ^2 \frak g_{-1}  \otimes \mathfrak g_{-1})^* \otimes U_{-1}.
    \end{eqnarray*}
\end{enumerate}

\end{proposition}

\begin{proof}
By Proposition \ref{prop: computation of H2 cohomology} and Lemma \ref{lem:H1 vanishing semisimple} and Lemma \ref{lem:H2 vanishing semisimple}, $H^{ 2}(\frak m, \frak g)_k$ vanishes for any $ k \geq 3$. By the uniqueness of $\zeta$ in the expression $\phi = \partial \eta + \zeta$ in Proposition \ref{prop: computation of H2 cohomology}, the nonvanishing $H^2(\frak m, \frak g)_k$ can be thought of as a subspace of the spaces in the right hand side.  In the first case, we can regards $\zeta_{\frak l}({Y^{-1}}, \cdot ) \in \Hom(\frak l_{-1},\frak l_{-1})$ because $\dim \frak l_{-2} =1$ implies that $\zeta_{\frak l}({Y^{-1}}, \cdot )\in \Hom(\frak l_{-2},\frak l_{-2})$ is just a constant multiple which comes from a boundary.
\end{proof}


 \subsection{Technical Lemmata} \label{sect:technical lemma}
%
We denote the restriction of $\partial$ to $\frak l_-$ by $\partial_0$, so that we have a subcomplex
 \begin{eqnarray*}
 0 \stackrel{\partial_0}{\longrightarrow} \frak g \stackrel{\partial_0}{\longrightarrow} \Hom( \frak l_-, \frak g ) \stackrel{\partial_0}{\longrightarrow} \Hom(\wedge^2 \frak l_-  ,\frak g   ) \stackrel{\partial_0}{\longrightarrow} \dots
 \end{eqnarray*}

 Recall that in the proof of the vanishing of $H^1(\frak m, \frak g)_k$ for positive $k$ (Proposition 48 of \cite{Ki}) a main difficulty is to show that  the nonvanishing  cohomology $H^1(\frak l_-, \mathfrak z)_1$ does not contribute to the cohomology $H^1(\frak m, \frak g)_1$, and a crucial Lemma is   the following.

\begin{lemma} [Lemma 27 (3) of \cite{Ki}] \label{lem: cohomology lemma 1 old version} For  $A \in U_0$, if the image of $\partial_0 A: \mathfrak l_{-1} \rightarrow U_{-1}$ has dimension $\leq 1$, then we have $\partial_0A=0$.
\end{lemma}

Under the assumption that $X$ is either of type $(B_m, \alpha_{m-1}, \alpha_m)$   or of type $(F_4, \alpha_2, \alpha_3)$, Lemma \ref{lem: cohomology lemma 1 old version}  can be improved.

\begin{lemma}   \label{lem: cohomology lemma 1} For  $A \in U_0$, if the image of $\partial_0 A: \mathfrak l_{-1} \rightarrow U_{-1}$ has dimension $\leq     2$, then we have $\partial_0A=0$.
\end{lemma}
%
Similarly, we have the following Lemma.

\begin{lemma} \label{lem: cohomology lemma 2} 
For $A \in \Hom(\frak l_-, U)_1$, if the image of $\partial_0A: \wedge^2 \frak l_{-1} \rightarrow U_{-1}$ has dimension $\leq 1$, then we have $\partial_0 A=0$.


\end{lemma}

\begin{proof}
Let $\{x_{-\alpha}\}$ be a basis of $\frak l_{-1}$ consisting of root vectors.  We may assume that $[x_{-\alpha}. x_{-\beta}]=c_{\alpha \beta}x_{-\alpha-\beta}$, $c_{\alpha \beta} \in \mathbb C$,  form a basis of $\frak l_{-2}$.
 Let $\{u_{\mu}\}$ ($\{u_{\lambda}\}$, respectively) be a basis of $U_{-1}$ ($U_0$, respectively) consisting of weight vectors. We may assume that $[x_{-\alpha}, u_{\lambda}] = u_{-\alpha+\lambda}$ if $-\alpha + \lambda$ is a weight.\\

    For $A \in \Hom(\frak l_{-}, U)_1$, we have
 \begin{eqnarray*}
 A(x_{-\alpha}) &=& \sum_{\lambda} A_{\lambda, \alpha} u_{\lambda} \\
 A(x_{-\alpha -\beta})&=& \sum_{\mu}A_{\mu, \alpha+\beta} u_{\mu},
 \end{eqnarray*}
 for  some matrices $A_{\lambda, \alpha}$ and $A_{\mu, \alpha+\beta}$.
 Then
 \begin{eqnarray*}
 (\partial_0A)(x_{-\alpha}, x_{-\beta}) &=&[x_{-\alpha}, A(x_{-\beta})] - [x_{-\beta}, A(x_{-\alpha}) - A([x_{-\alpha}, x_{-\beta}]) \\
&=& [x_{-\alpha}, \sum_{\lambda} A_{\lambda, \beta}u_{\lambda}] -[x_{-\beta}, \sum_{\lambda}A_{\lambda, \alpha}u_{\lambda}] -  A(x_{-\alpha-\beta}c_{\alpha \beta})  \\
&=& \sum_{\lambda} A_{\lambda, \beta} u_{-\alpha+\lambda} - \sum_{\lambda} A_{\lambda, \alpha} u_{-\beta+\lambda}- \sum_{\mu} A_{\mu, \alpha+\beta}c_{\alpha \beta} u_{\mu} \\
&=&\sum_{\mu}  \left( A_{\mu+ \alpha, \beta} - A_{\mu+\beta, \alpha} + A_{\mu, \alpha+\beta} c_{\alpha \beta}  \right) u_{\mu}.
 \end{eqnarray*}
 %
The condition that the image of $\partial_0A $ has dimension $\leq 1$ implies that for any choice of a pair $(x_{-\alpha}, x_{-\beta})$, $(\partial_0A)(x_{-\alpha}, x_{-\beta})$ is parallel to each other. \\

 \begin{enumerate}
 \item If $(\frak m, \frak g_0)$ is of type $(B_m, \alpha_{m-1}, \alpha_m)$, the action $ \frak l_1 \times U_{-1} \rightarrow U_0$ is given by
 \begin{eqnarray*}
%
 (\Sym ^2 V^* \otimes W^*) \times  V  &\rightarrow&  V^* \otimes W^* \\
%
(v_i^{*2} \otimes w^*, v_j) &\mapsto& \delta_{ij}v_i \otimes w^*
 \end{eqnarray*}
where $\{v_i : i=1,2\}$ is a basis of $V$ and $w$ is an element of $W$. Furthermore,  for $i=1,2$, the Lie bracket of any two element of $\{v_i^{*2} \otimes w^*: w \in W\}$ is zero. Write $v_1, v_2$ as $u_{\mu}, u_{\nu}$. Then
  $(\partial_0A)(x_{-\alpha}, x_{-\beta})$ is given by  $$\left( A_{\mu+ \alpha, \beta} - A_{\mu+\beta, \alpha} + A_{\mu, \alpha+\beta} c_{\alpha \beta}  \right) u_{\mu} + \left( A_{\nu+ \alpha, \beta} - A_{\nu+\beta, \alpha} + A_{\nu, \alpha+\beta} c_{\alpha \beta}  \right) u_{\nu}. $$
 If we take $x_{-\alpha}, x_{-\beta}$ in $\{v_1^2 \otimes w : w \in W\}$, then the coefficient of $u_{\mu}$ is zero. If we take $x_{-\alpha}, x_{-\beta}$ in $\{v_2^2 \otimes w : w \in W\}$, then the coefficient of $u_{\nu}$ is also zero.
 Since  $(\partial_0A)(x_{-\alpha}, x_{-\beta})$ is parallel to each other for any choice of a pair $(x_{-\alpha},x_{-\beta})$, both coefficients are zero for any $\alpha, \beta$. \\

 \item If $(\frak m, \frak g_0)$ is of type $(F_4, \alpha_2, \alpha_3)$, the action $ \frak l_1 \times U_{-1} \rightarrow U_0$ is given by
 \begin{eqnarray*}
 (\Sym^2 V^* \otimes W^*) \times  V &\rightarrow& V^* \otimes W^*\\
 ( {v_i^*}\circ{v_j^*} \otimes w^* , v_k ) &\mapsto & (\delta_{ik}v_j^* + \delta_{jk}v_i^*) \otimes w^*.
  \end{eqnarray*}
In particular, $({v_i^*}^2 \otimes w^*, v_j )$ maps to $\delta_{ij}v_i^* \otimes w^*$, where $\{v_i : i=1,2,3\}$ is a basis of $V$ and $w$ is an element of $W$.  Furthermore, for  $i=1,2,3$,  the Lie bracket of any two elements in $\{ {v_i^*}^2 \otimes w: w \in W\}$ is zero.

Write $v_1, v_2,v_3$ as $u_{\mu}, u_{\nu}, u_{\xi}$. Then
  $(\partial_0A)(x_{-\alpha}, x_{-\beta})$ is given by
  \begin{eqnarray*}
 && \left( A_{\mu+ \alpha, \beta} - A_{\mu+\beta, \alpha} + A_{\mu, \alpha+\beta} c_{\alpha \beta}  \right) u_{\mu} + \left( A_{\nu+ \alpha, \beta} - A_{\nu+\beta, \alpha} + A_{\nu, \alpha+\beta} c_{\alpha \beta}  \right) u_{\nu} \\
 && + \left( A_{\xi+ \alpha, \beta} - A_{\xi+\beta, \alpha} + A_{\xi, \alpha+\beta} c_{\alpha \beta}  \right) u_{\xi}
 \end{eqnarray*}
 By the same arguments as in (1), we get that all coefficients are zero. 
 \end{enumerate}
\end{proof}


\subsection{A proof of Proposition \ref{prop: computation of H2 cohomology}}   \label{sect: proof of proposition}
%
We list up properties of $\frak g = U + \frak z + \frak l$.

\begin{enumerate}
\item [(${\bf P_0}$)] $[\frak l, \frak z]=0$ and $[\frak l_-, U_-]=0$
\item [(${\bf P_1}$)] For $X \in \frak l_{\geq - \mu+1} +U_{\geq 0} $, if $[\frak l_-, X]=0$, then $X=0$. \\
\end{enumerate}

\begin{proof} [Proof of Proposition \ref{prop: computation of H2 cohomology}]
 Let $\phi \in \Hom(\wedge^2 \frak m, \frak g)_k$ be such that $\partial \phi=0$, where $k \geq 1$. We will show that there exist   $\eta \in \Hom(\frak m, \frak g)_k$, $\zeta=\zeta_{\frak l}+\zeta_U$ with $\zeta_{\frak l} \in  H^2(\frak l_-, U)_k$ and $\zeta_U \in H^2(U_-, U)_1$, and $\xi=\xi_{\frak l} + \xi_U $ with $\xi_U(X^{-1}, \,\cdot\,) \in H^1(\frak l_-, U)_{k-1}$ and $\xi_{\frak l}(X^{-1}, \,\cdot\,) \in H^1(\frak l_-, \frak l)_{k-1}$ for $X^{-1} \in U_-$ satisfying that  $$\phi = \partial \eta + \zeta +\xi .$$ In other words, for any $X^{-1}, Y^{-1} \in U_-$ and $X^0, Y^0 \in \frak l_-$, we have
 \begin{eqnarray*}
 \phi(X^{-1} + X^{0}, Y^{-1}+ Y^0)   &=&  [X^{-1}+X^0, \eta(Y^{-1}+ Y^0)] -[Y^{-1}+Y^0, \eta(X^{-1} + X^0)]\\
 &&- \eta([X^0, Y^0])  + \zeta(X^0, Y^0) +\xi(X^0, Y^{-1})\\ && + \xi(X^{-1}, Y^0) + \zeta(X^{-1}, Y^{-1}),
 \end{eqnarray*}
 or equivalently,
 \begin{eqnarray*}
 \phi_{\frak l+\frak z}(X^{-1} + X^0, Y^{-1}+Y^0)
   &=& [X^0, \eta_{\frak l+\frak z}(Y^{-1} +Y^0)] - [Y^0, \eta_{\frak l+\frak z}(X^{-1} + X^0)] \\
  &&- \eta_{\frak l+\frak z}([X^0, Y^0])+\xi_{\frak l}(X^0, Y^{-1}) + \xi_{\frak l}(X^{-1}, Y^0) \\
   \phi_U(X^{-1}+X^0, Y^{-1}+Y^0)
 &=& [X^{-1}, \eta_{\frak l + \frak z}(Y^{-1} +Y^0)]-[Y^{-1}, \eta_{\frak l + \frak z}(X^{-1}+X^0)] \\
 &&  + [X^0, \eta_U(Y^{-1} + Y^0)] -[Y^0, \eta_U(X^{-1} + X^0)]   \\
 && - \eta_U([X^0, Y^0]) + \zeta_U(X^0, Y^0)+\xi_U(X^0, Y^{-1}) \\ && + \xi_U(X^{-1}, Y^0) + \zeta_U(X^{-1}, Y^{-1}).
 \end{eqnarray*}

\vskip 10 pt
 To do this, we decompose the identity
   $$\partial \phi(X^{-1}+ X^0, Y^{-1} + Y^0, Z^{-1} +Z^0)=0$$ into the sum of two identities: \\

%{\small \vskip 10 pt 

%\begin{itemize}
%\item[{(\rm I)}]
\noindent (I) $[X^0, \phi_{\frak l+ \frak z}(Y^{-1}+Y^0, Z^{-1}+Z^0)] -[Y^0, \phi_{\frak l+ \frak z}(X^{-1}+X^0 , Z^{-1} +Z^0)]
+ [Z^0, \phi_{\frak l+\frak z}(X^{-1}+ X^0, Y^{-1} +Y^0)] -\phi_{\frak l+ \frak z}([X^0, Y^0],Z^{-1}+ Z^0) + \phi_{\frak l + \frak z}([X^0,Z^0],Y^{-1}+Y^0) - \phi_{\frak l + \frak z}([Y^0, Z^0], X^{-1} +X^0) =0$. \vskip 10 pt

%\item [{(\rm II)}]
\noindent (II) $[X^0, \phi_U(Y^{-1}+Y^0,Z^{-1}+Z^0)] - [Y^0, \phi_U(X^{-1}+X^0, Z^{-1}+Z^0)] +[Z^0, \phi_U(X^{-1}+X^0, Y^{-1}+Y^0)] +[X^{-1}, \phi_{\frak l + \frak z}(Y^{-1}+Y^0, Z^{-1}+Z^0)] -[Y^{-1}, \phi_{\frak l +\frak z}(X^{-1}+X^0, Z^{-1}+Z^0)] +[Z^{-1}, \phi_{\frak l + \frak z}(X^{-1}+X^0, Y^{-1}+Y^0)]  -\phi_U([X^0, Y^0],Z^{-1}+Z^0) + \phi_U([X^0, Z^0], Y^{-1}+Y0) - \phi_U([Y^0, Z^0], X^{-1}+X^0)=0$.
%\end{itemize}
\vskip 10 pt
%}

 \noindent {\bf  R1.} There exist    $\eta_{\frak l+\frak z}  +\eta_{U} \in \Hom(\frak l_-, \frak g)_k $ and $\zeta_{\frak l+ \frak z}  + \zeta_U \in H^2(\frak l_-, \frak g)_k$ such that
 \begin{eqnarray*}
 \phi_{\frak l+\frak z}(X^0, Y^0) &= & [X^0, \eta_{\frak l+ \frak z}(Y^0)] -[Y^0, \eta_{\frak l+\frak z}(X^0)] - \eta_{\frak l+\frak z}([X^0, Y^0]) + \zeta_{\frak l+\frak z}(X^0, Y^0)   \\
 %\phi_{\frak z}(X^0, Y^0) &=& \hskip 4.7 cm  -\psi_{\frak z}([X^0,Y^0]) + \zeta_{\frak z}(X^0,Y^0) \\
 \phi_{U}(X^0,Y^0) &=&  [X^0, \eta_U(Y^0)] -[Y^0, \eta_U(X^0)] -\eta_U([X^0,Y^0])+ \zeta_{U}(X^0, Y^0)
 \end{eqnarray*}
 %
 Indeed, putting $X^{-1}=Y^{-1}=Z^{-1}=0$ into (I) and (II), we get $\partial_0(\phi|_{\frak l_-})=0$.

Note that    $\zeta_{\frak l+ \frak z} = \zeta_{\frak z}$ if $k=2$, and  $\zeta_{\frak l+ \frak z}=0$, otherwise,  by Lemma \ref{lem:H2 vanishing semisimple} (i) and (ii).  \\

\noindent {\bf  R2.} We have  $\phi_{\frak l_{r}}(Y^{-1}, Z^{-1})=0$ for $r \geq - \mu +1$ and $r \neq 0$. Indeed,
 putting $Y^0=Z^0=0$ to (I) for $k\neq 2$, we get $[X^0, \phi_{\frak l_r}(Y^{-1},Z^{-1}) ]=0$ for $r\neq 0$. The desired vanishing follows from (${\bf P_1}$).
 \\

  \noindent {\bf  R3.} We may extend $\eta_{\frak l + \frak z}$ from $\frak l_-$ to $\frak m=\frak l_- + U_-$ %$\eta_{\frak l+\frak z} \in \Hom(\frak l_-+U_{-}, \frak l+ \frak z)$ and
   so that
  $$\phi_{\frak l+\frak z}(Y^{-1}, \,\cdot\,) = [\eta_{\frak l+\frak z }(Y^{-1}),\,\cdot\,] -\xi_{\frak l+\frak z}^{Y^{-1}}(\,\cdot\,) \text{ in }\Hom(\frak l_-, \frak l+ \frak z)_{k-1}$$
  for some $\xi_{\frak l+\frak z}^{Y^{-1}} \in H^1(\frak l_-, \frak l+\frak z)_{k-1}$.
  Indeed, putting $X^{-1} =Y^0 =Z^{-1}=0$ to (I) we get
  $$[X^0, \phi_{\frak l+ \frak z}(Y^{-1},Z^0)] - [Z^0, \phi_{\frak l + \frak z}(Y^{-1}, X^0)] - \phi_{\frak l+ \frak z}(Y^{-1}, [X^0, Z^0]) =0,$$
  or equivalently,
  \begin{eqnarray*}
  \partial_0(\phi_{\frak l+\frak z}(Y^{-1}, \,\cdot\,))=0.  %\\
  \end{eqnarray*}
 Thus
  there exit $\eta_{\frak l + \frak z}^{Y^{-1}} \in (\frak l + \frak z)_{k-1}$ and   $\xi_{\frak l+\frak z}^{Y^{-1}}(\,\cdot\,) \in H^1(\frak l_-, \frak l+\frak z)_{k-1}$ such that
 $$\phi_{\frak l+\frak z}(Y^{-1}, \,\cdot\,) = [\eta_{\frak l+\frak z }^{Y^{-1}},\,\cdot\,] -\xi_{\frak l+\frak z}^{Y^{-1}}(\,\cdot\,) .$$
  Note that $\xi_{\frak l+ \frak z}^{Y^{-1}} =\xi_{\frak l}^{Y^{-1}}$ if $k=1$ (for $B_3$ type), $\xi_{\frak l+ \frak z}^{Y^{-1}} =\xi_{\frak z}^{Y^{-1}}$ if $k=2$ and $\xi_{\frak l+\frak z}^{Y^{-1}}=0$ otherwise, by Lemma \ref{lem:H1 vanishing semisimple} (i) and (ii). %(${\bf P_2}$)
Define $\eta_{\frak l + \frak z}(Y^{-1})$ by $\eta_{\frak l + \frak z}(Y^{-1}):= \eta_{\frak l+ \frak z}^{Y^{-1}}$.     \\

\noindent {\bf R4.} We have  $\phi_{\frak l_0 + \frak z}(Y^{-1}, Z^{-1}) =0 $. Indeed,  putting $Y^0=Z^0=0$ to (I) for $k=2$, we get $[X^0, \phi_{\frak l_0 +\frak z}(Y^{-1},Z^{-1}) ]=0$ and putting $X^0=Y^0=Z^0=0$ to (II) for $k=2$ we get
  \begin{eqnarray*}
  [X^{-1}, \phi_{\frak l_0 + \frak z}(Y^{-1}, Z^{-1})] - [Y^{-1}, \phi_{\frak l_0+ \frak z}(X^{-1}, Z^{-1})] + [Z^{-1}, \phi_{\frak l_0+ \frak z}(X^{-1}, Y^{-1})] =0.
  \end{eqnarray*}
  Take independent vectors  $X^{-1}, Y^{-1}, Z^{-1}$  in $U_-$ and arbitrary $X^0\in \frak l_-$ to get the desired result. \\

  \noindent {\bf R5.} We may extend $\zeta_U$ from $\frak l_-$ to $\frak m = \frak l_-+ U_-$ such that $\zeta|_{\wedge^2U_-} \in H^2(U_-, U)_1$, satisfying
  $$\phi_U(Y^{-1}, Z^{-1}) =[Y^{-1}, \eta_{\frak l+ \frak z}(Z^{-1})] - [Z^{-1}, \eta_{\frak l + \frak z}(Y^{-1})] +\zeta_U(Y^{-1}, Z^{-1})$$
   and
   $$\phi_{\frak z}(X^0, Z^{-1})% =[X^0, \eta_{\frak z}(Z^{-1})]+\xi_{\frak z}^{Z^{-1}}(X^0)
   = 0.$$

Indeed, putting $X^{-1}=Y^0=Z^0=0$ to (II) we get
  \begin{eqnarray*}
  [X^0, \phi_U(Y^{-1}, Z^{-1})]
     - [Y^{-1}, \phi_{\frak l+ \frak z}( X^{0}, Z^{-1})] + [Z^{-1}, \phi_{\frak l+ \frak z}( X^0, Y^{-1})] =0.
  \end{eqnarray*}
Thus
  \begin{eqnarray*}
  &&[X^0, \phi_U(Y^{-1}, Z^{-1})] \\
   &\stackrel{\bf R3.}{=}& [Y^{-1},[X^0, \eta_{\frak l + \frak z}(Z^{-1})] +\xi_{\frak l+\frak z}^{Z^{-1}}(X^0)] - [Z^{-1},[X ^0, \eta_{\frak l + \frak z}(Y^{-1})] +\xi_{\frak l+\frak z}^{Y^{-1}}(X^0)]  \\
   &\stackrel{(\bf P_0)}{=}&  [X^0,[Y^{-1}, \eta_{\frak l + \frak z}(Z^{-1})]]  - [X^0,[ Z^{-1}, \eta_{\frak l + \frak z}(Y^{-1})]]
     +[Y^{-1}, \xi_{\frak l+\frak z}^{Z^{-1}}(X^0)]  -[Z^{-1}, \xi_{\frak l+\frak z}^{Y^{-1}}(X^0)].
  \end{eqnarray*}
  %By (${\bf P_0}$)
  Hence
  \begin{eqnarray*}
   &&[X^0, \phi_U(Y^{-1}, Z^{-1})]
   -[X^0,[Y^{-1}, \eta_{\frak l + \frak z}(Z^{-1})]]  + [X^0,[ Z^{-1}, \eta_{\frak l + \frak z}(Y^{-1})]] \\
   &&=  [Y^{-1}, \xi_{\frak l+\frak z}^{Z^{-1}}(X^0)]  -[Z^{-1}, \xi_{\frak l+\frak z}^{Y^{-1}}(X^0)].
    \end{eqnarray*}
    Therefore, we have
    \begin{eqnarray*}
    \partial_0\left( \phi_U(Y^{-1}, Z^{-1}) -[Y^{-1}, \eta_{\frak l+ \frak z}(Z^{-1})] + [Z^{-1}, \eta_{\frak l + \frak z}(Y^{-1})]\right)
     =[Y^{-1}, \xi_{ \frak l+\frak z}^{Z^{-1}}(\,\cdot\,)]  -[Z^{-1}, \xi_{\frak l+ \frak z}^{Y^{-1}}(\,\cdot\,)].
    \end{eqnarray*}
    We claim that both sides vanish.

   By Lemma \ref{lem:H1 vanishing semisimple} (i) and (ii), $\xi_{\frak l+ \frak z}^{Y^{-1}}$ is $\xi_{\frak l}^{Y^{-1}}$ or zero if $k=1$, $\xi_{\frak l+ \frak z}^{Y^{-1}} $ is  $\xi_{\frak z}^{Y^{-1}}$ if $k=2$, and is zero otherwise.
   In the last case,
   $$\partial_0\left( \phi_U(Y^{-1}, Z^{-1}) -[Y^{-1}, \eta_{\frak l+ \frak z}(Z^{-1})] + [Z^{-1}, \eta_{\frak l + \frak z}(Y^{-1})]\right) =0.$$

   In the second case,
   $$A^{Y^{-1},Z^{-1}}:=\phi_U(Y^{-1}, Z^{-1}) -[Y^{-1}, \eta_{\frak l+ \frak z}(Z^{-1})] + [Z^{-1}, \eta_{\frak l + \frak z}(Y^{-1})]$$
   is an element of $U_0$ and we have
    $$\partial_0 \,A^{Y^{-1},Z^{-1}} =[Y^{-1}, \xi_{ \frak z}^{Z^{-1}}(\,\cdot\,)]  -[Z^{-1}, \xi_{ \frak z}^{Y^{-1}}(\,\cdot\,)].$$
   Both sides vanish by Lemma \ref{lem: cohomology lemma 1} because $[Y^{-1},\mathfrak z] -[Z^{-1}, \mathfrak z]$ has dimension two.
    In the first case, it is clear that both sides vanish because $[U_, \, \frak g_-]=0$.

   Therefore, we have
   $$\phi_U(Y^{-1}, Z^{-1}) =[Y^{-1}, \eta_{\frak l+ \frak z}(Z^{-1})] - [Z^{-1}, \eta_{\frak l + \frak z}(Y^{-1})]+\zeta_U(Y^{-1}, Z^{-1}) \text{ by } {\rm(}{\bf P_1}{\rm)} $$
   and
   $$\phi_{\frak z}(X^0, Z^{-1}) = \xi_{\frak z}^{Z^{-1}}(X^0) = 0.$$
 Note that $\zeta_U(Y^{-1}, Z^{-1}) \in U_-$.\\

%
   \noindent  {\bf R6.} We may extend $\eta_U$ from $\frak l_-$ to $\frak m=\frak l_- + U_-$ so that   %$\eta_U^{Y^{-1}} \in U$ such that
  $$\phi_U(Y^{-1},\,\cdot\,) - [Y^{-1}, \eta_{\frak l+ \frak z}(\,\cdot\,)] =-[\,\cdot\,, \eta_U(Y^{-1})] \text{ in } \Hom(\frak l_-, U)_{k-1}, $$
   and $$\zeta_{\frak z} (X^0,\,Y^0)=0.$$
 %
  Indeed, putting $X^{-1}=Y^0=Z^{-1}=0$ we get
   \begin{eqnarray*}
   &&[X^0, \phi_U(Y^{-1}, Z^0) + [Z^0, \phi_U(X^0, Y^{-1})] -[Y^{-1}, \phi_{\frak l+ \frak z}(X^0), Z^0)] \\
   && + \phi_U([X^0, Z^0], Y^{-1}) =0.
   \end{eqnarray*}
   %By {\bf R1.}
   Thus
   \begin{eqnarray*}
   &&[X^0, \phi_U(Y^{-1},Z^0)] + [Z^0, \phi_U(X^0, Y^{-1})] + \phi_U([X^0, Z^0],Y^{-1}) \\
   &=& [Y^{-1}, \phi_{\frak l + \frak z}(X^0, Z^0)] \\
   &\stackrel{\bf R1.}{=}&[Y^{-1}, [X^0, \eta_{\frak l + \frak z}(Z^0)] - [Z^0, \eta_{\frak l + \frak z}(X^0)]-\eta_{\frak l+ \frak z}([X^0, Z^0]) + \zeta_{\frak l+ \frak z}(X^0, Z^0) ] \\
   &\stackrel{({\bf P_0})}{=}& [X^0, [Y^{-1}, \eta_{\frak l + \frak z}(Y^{0})]] - [Y^{-1}, [Z^0, \eta_{\frak l + \frak z}(Y^0)]] -[Y^{-1}, \eta_{\frak l + \frak z}([X^0, Z^0])] \\
   && + [Y^{-1}, \zeta_{\frak l+ \frak z}(X^0, Z^0)  ]
   \end{eqnarray*}
   Therefore, we have
   \begin{eqnarray*}
   \partial_0(\phi_U(Y^{-1},\,\cdot\,) -[Y^{-1}, \eta_{\frak l + \frak z}(\,\cdot\,)] ) (\,\cdot\,,\,\cdot\,) =[Y^{-1}, \zeta_{\frak l+ \frak z}(\,\cdot\,,\,\cdot\,)] .
   \end{eqnarray*}
    We claim that both  sides vanish.


   By Lemma \ref{lem:H2 vanishing semisimple} (i) and (ii),   $\zeta_{\frak l+ \frak z}$  is $\zeta_{\frak z} $ if $k=2$, and is zero, otherwise.
   In the second case,
   $$\partial_0(\phi_U(Y^{-1},\,\cdot\,) -[Y^{-1}, \eta_{\frak l + \frak z}(\,\cdot\,)] ) (\,\cdot\,,\,\cdot\,) = 0.$$
    In the first case,  $$A^{Y^{-1}}:=\phi_U(Y^{-1},\,\cdot\,) -[Y^{-1}, \eta_{\frak l + \frak z}(\,\cdot\,)] $$ is an element of $\Hom(\frak l_-, U)_1$, and
   we have
    \begin{eqnarray*}
   \partial_0 A^{Y^{-1}} =[Y^{-1}, \zeta_{\frak z}(\,\cdot\,,\,\cdot\,)].
   \end{eqnarray*}
    %
   Both sides vanish by Lemma \ref{lem: cohomology lemma 2} because $[Y^{-1},\frak z]$ has dimension one.

   By Lemma \ref{lem:H1 vanishing semisimple} %(${\bf P_2}$),
   there exist    $\eta_U^{Y^{-1}} \in U_{k-1}$ and $\xi_U^{Y^{-1}} \in H^1(\frak l_-, U)_{k-1}$ such that
  $$\phi_U(Y^{-1},\,\cdot\,) - [Y^{-1}, \eta_{\frak l+ \frak z}(\,\cdot\,)] =-[\,\cdot\,, \eta_U^{Y^{-1}}] +\xi_U^{Y^{-1}}(\,\cdot\,) \text{ in } \Hom(\frak l_-, U)_{k-1}, $$
  and $$\zeta_{\frak z} (\,\cdot\,,\,\cdot\,)=0.$$
  Define $\eta_U(Y^{-1})$ by $\eta_U(Y^{-1}):=\eta_U^{Y^{-1}}$ and define $\xi_U \in \Hom (\frak l_{-1} \wedge U_{-1}, U_{-1})$ by $\xi_U(Y^{-1}, X^0)=-\xi_U(X^0, Y^{-1}):=\xi_U^{Y^{-1}}(X^0)$. \\


Consequently, there exist $\eta = \eta_{\frak l+ \frak z} + \eta_U \in \Hom(\frak m, \frak g)_k$, $\zeta =  \zeta_U \in H^2(\frak l_-, U)_k+H^2(U_-, U)_1$ and $\xi_U+ \xi_{\frak l} \in (\frak l_{-1} \wedge U_{-1})^* \otimes (U_{-1}+ \frak l_{-1})$ satisfying
\begin{eqnarray*}
\phi_{\frak l+\frak z}(X^0, Y^0) &\stackrel{{\bf R1.}, {\bf R6}}{=}& [X^0, \eta_{\frak l+\frak z}(Y^0)] -[Y^0,\eta_{\frak l+\frak z}(X^0)] - \eta_{\frak l+\frak z}([X^0,Y^0])\\% + \zeta_{\frak l+ \frak z}(X^0,Y^0)  \\
\phi_{\frak l+\frak z}(X^0, Y^{-1})&\stackrel{{\bf R3.},{\bf R5.}}{=}&[X^0, \eta_{\frak l}(Y^{-1})]+\xi_{\frak l}(X^0, Y^{-1}) \\
  \phi_{\frak l_{\geq -\mu+1}+\frak z}(X^{-1},Y^{-1}) &\stackrel{{\bf R2.}, {\bf R4.}}{=}&0 \\
  \phi_U(X^0, Y^0)&\stackrel{\bf R1.}{=}& [X^0, \eta_U(Y^0)] -[Y^0,\eta_U(X^0)] - \eta_U([X^0, Y^0]) + \zeta_U(X^0,Y^0) \\
  \phi_U(X^0, Y^{-1}) &\stackrel{\bf R6.}{=}& -[Y^{-1}, \eta_{\frak l+\frak z}(X^0)] +[X^0, \eta_U(Y^{-1})] +\xi_U (X^0, Y^{-1}) \\
  \phi_U(X^{-1}, Y^{-1}) &\stackrel{\bf R5.}{=}& [X^{-1}, \eta_{\frak l+ \frak z}(Y^{-1})] -[Y^{-1},\eta_{\frak l+\frak z}(X^{-1})]+\zeta_U(Y^{-1}, Z^{-1}).
\end{eqnarray*}
%

  It remains to show that the choice of $\zeta$ in the  expression $\phi = \partial \eta + \zeta   $ is unique. Indeed, if $\partial \eta + \zeta   =\partial \eta' + \zeta'   $, then $\partial(\eta-\eta') = \zeta' - \zeta$. It suffices to show that if $\partial \eta = \zeta$, then   $\zeta$ is zero.
From
$$(\partial \eta)(X^{-1}+ X^0, Y^{-1}+ Y^0) = \zeta(X^{-1}+X^0, Y^{-1}+Y^0)$$
it follows that
\begin{eqnarray*}
\zeta(X^0, Y^0) &=&[X^0, \eta(Y^0)]-[Y^0, \eta(X^0)] -\eta([X^0, Y^0]) \\
\zeta(X^0, Y^{-1})&=& [X^0, \eta(Y^{-1})] - [Y^{-1}, \eta(X^0)].
\end{eqnarray*}
The first identity implies that $\zeta = \partial_0 \eta|_{\frak l_-}$. Since $\zeta|_{\wedge^2 \frak l_-} \in H^2(\frak l_-, U )_k$, $\zeta|_{\wedge^2 \frak l_-}$ vanishes and so does $\partial_0 \eta|_{\frak l_-}$. Since $H^1(\frak l_-, \frak g)_k=0$ for any $k\geq 1$, there is $\chi$ such that $\eta|_{\frak l_-} =\partial_0\chi$. Then $\eta(X^0) =[\chi, X^0]$ for $X^0 \in \frak l_-$.

The second identity becomes
\begin{eqnarray*}
\zeta(X^0, Y^{-1}) &=& [X^0, \eta(Y^{-1})] -[Y^{-1}, [\chi, X^0]] \\
&=& [X^0, \eta(Y^{-1}) + [Y^{-1}, \chi]] \text{ because } [\frak l_-,   U_-]=0
\end{eqnarray*}
Since $\zeta(\,\cdot\,, Y^{-1}) |_{\frak l_-} \in H^1(\frak l_-, U)_{k-1}$, we have  $\zeta(\,\cdot\,, Y^{-1}) |_{\frak l_-} =0$.


This completes the proof of Proposition \ref{prop: computation of H2 cohomology}.
  \end{proof}



\section{Local equivalence of geometric structures} \label{sect:local equivalence}



We keep the same notations as in Section \ref{sect:horospherical BF type}.

\begin{proposition} \label{extension to rationally connected open subset}
Let $M$ be a Fano manifold of Picard number one and $\mathcal C_x(M) \subset \mathbb PT_xM$ be the variety of minimal rational tangents at a general point $x \in X$ associated to a minimal dominant rational component $\mathcal K$. Suppose that $\mathcal C_x(M) \subset \mathbb PT_xM$ is projectively equivalent to ${\bf S} \subset \mathbb P \frak m$ for general $x \in M$. Denote by $D \subset TM$   the distribution on $M$ obtained by the  linear span of the affine cone of $\mathcal C_x(M)$ in $T_xM$.
Then there is a Zariski open subset $M^0\subset M$ such that a general member  of $\mathcal K$ lies on $M^0$, satisfying  the following properties:
\begin{enumerate}
\item $D|_{M^0}$ is of type $\frak m$;
\item $\mathcal C|_{M^0}\subset \mathbb P D|_{M^0}$ defines an ${\bf S}$-structure on $(M^0, D|_{M^0})$ and there corresponds to a $G_0$-structure on $(M^0, D|_{M^0})$, where $G_0=G(\widehat{\bf S})$.
\end{enumerate}

\end{proposition}


\begin{proof}
The distribution $D$ is holomorphic outside its singular set $Sing(D) \subset M$ of codimension $\geq 2$.
 Let $M' \subset M-Sing(D)$ be a  Zariski open subset such that for all $x \in M'$ satisfying $\mathcal C_x \subset \mathbb P D_x$ is projectively equivalent to $ {\bf S} \subset \mathbb P {\frak g_{-1}}$.

Let $C $ be a standard minimal rational curve represented by $f:\mathbb P^1 \rightarrow M$ with $C \cap M' \not = \emptyset$ and $C \not\subset bad(\mathcal K)$ and  $C \cap Sing(D) =\emptyset$. Then for a generic point $y \in C$, $\mathcal C_y \subset \mathbb PD_y$ is projectively equivalent to $\bf S\subset \mathbb P \frak g_{-1}$. By Proposition \ref{parallel transport of II} and Proposition \ref{parallel transport of III} and Lemma  \ref{parallel transport of III-X case}, the relative second fundamental form and the relative third fundamental form of $\mathcal C(M)$ along the lifting $C^{\sharp}$ of $C$ is constant, and we have %the following.
\begin{eqnarray*}
f^*D &=& \mathcal O(2) \oplus \mathcal O(1)^p \oplus \mathcal O^r \oplus \mathcal O(-1)^s \\
f^*(TM/D) &=& \mathcal O(1)^s \oplus \mathcal O^t.
\end{eqnarray*}
For the values $p,q,r,s,t$, see Remark \ref{pqrst}.
Furthermore,
the pull-back $(f^{\sharp})^*\widehat{T}^{\varpi} $ of the relative affine tangent bundle $ \widehat{T}^{\varpi} $ of $\mathcal C \subset \mathbb P(TM)$ is the positive part $ P:=\mathcal O(2) \oplus \mathcal O(1)^p$ of $f^*D$, and
   the pull-back $(f^{\sharp})^*\widehat{T}^{(2), \varpi} $ of the relative  second osculating affine   bundle   $ \widehat{T}^{(2), \varpi} $ of $\mathcal C \subset \mathbb P(TM)$  is the subbundle  $P^{(2)} :=\mathcal O(2) \oplus \mathcal O(1)^p \oplus \mathcal O^r$ of $f^*D$. % where $r$ is the dimension of the image of $II^{\varpi}$.

 By Proposition \ref{parallel transport of vmrt}, for any $y\in C$, $\mathcal C_y \subset \mathbb PD_y$ is projectively equivalent to $\bf S\subset \mathbb P \frak g_{-1}$.

Let $D^{-1} :=D$ and $D^{-2} := D + [D,D]$. For generic $x \in M$, the Frobenius bracket $ [\,,\,]$ is given by
\begin{eqnarray*}
    [\,,\,]_1&:&\wedge^2 D_x^{-1} \rightarrow D_x^{-2}/D_x \subset T_xM/D_x \\ [1 pt]
    [\,,\,]_2&:&D_x^{-1} \wedge D_x^{-2} \rightarrow (T_xM/D_x)/(D_x^{-2}/D_x)= T_xM/D_x^{-2}.
    \end{eqnarray*}
In particular, if $X$ is $(B_m, \alpha_{m-1}, \alpha_m)$, then $D^{-2}=TM$ and simply the bracket is
    \begin{eqnarray*} [\,,\,]:\wedge^2 D_x \rightarrow T_xM/D_x.
    \end{eqnarray*}

For generic $v \in \widehat {\mathcal C}_x$ and $w \in T_v(\widehat{\mathcal C}_x)$, we have $[v,w]=0$ by Proposition \ref{kernel of frobenius braket}. Hence, by Proposition \ref{fundamental graded Lie algebra of vmrt} (1) and its proof, the symbol algebra $\Symb_x(D)$ is isomorphic to $\mathfrak m$ at general $x$. It remains to show that $\Symb_y(D)$ is isomorphic to $\mathfrak m$ for any $y \in C$. For this, we need to know the decompositions of vector bundles related to the Frobenius bracket. From now on,  we use notations $D|_C$, $TM/D|_C$, etc., instead of $f^*D$, $f^*(TM/D)$, for simplicity.

\begin{enumerate}
\item  When $X$ is $(B_m, \alpha_{m-1}, \alpha_m)$,
\begin{enumerate}
\item $\dim \frak g_{-1} = 3m-1$, $\dim \widehat{\bf S} =m+1$, $\dim \frak g_{-2} = (m-1)(m-2)/2$
\item $D|_C=\mathcal O(2) \oplus \mathcal O(1)^m \oplus \mathcal O^m \oplus \mathcal O(-1)^{m-2}$, $TM/D|_C = \mathcal O(1)^{m-2} \oplus \mathcal O^t$, where $t=(m-2)(m-3)/2$.
\end{enumerate} %where $t=q-r-2s$

\item When $X$ is $(F_4, \alpha_2, \alpha_3)$,
\begin{enumerate}
\item $\dim \frak g_{-1} = 15$, $\dim \widehat{\bf S} =$, $\dim \frak g_{-2} = 6$, $\dim \frak g_{-3} = 2$
\item $D|_C=\mathcal O(2) \oplus \mathcal O(1)^4 \oplus \mathcal O^7 \oplus \mathcal O(-1)^{3}$, $TM/D|_C = \mathcal O(1)^{3} \oplus \mathcal O^5$.
\end{enumerate}

\end{enumerate}


\begin{lemma}   \label{lem:splitting type of V}
Let $\mathcal V$ and $\mathcal W $ be vector bundles on $M'$ associated with $V$ and $W$.
Let $C$ be a general member of $\mathcal K$ passing through $x \in M'$.
\begin{enumerate}
 \item  When $X$ is $(B_m, \alpha_{m-1}, \alpha_m)$, we have
$$\mathcal V|_C= \mathcal O(1) \oplus \mathcal O \text{ and } \mathcal W|_C= \mathcal O \oplus \mathcal O(-1)^{m-2}.  $$

 \item  When $X$ is $(F_4, \alpha_2, \alpha_3)$,
 we have
$$\mathcal V|_C= \mathcal O(1) \oplus \mathcal O^2 \text{ and } \mathcal W|_C= \mathcal O \oplus \mathcal O(-1)  .$$

\end{enumerate}

\end{lemma}

\begin{proof}
We adapt an argument in the proof of Proposition 7.3 of \cite{HL21}.

Recall that we have the following exact sequence:
\begin{eqnarray*}
0 \rightarrow \mathcal V|_C \rightarrow D |_C = \mathcal O(2) \oplus \mathcal O(1)^p \oplus \mathcal O^r \oplus \mathcal O(-1)^s \rightarrow \Sym^2 \mathcal V \otimes \mathcal W |_C \rightarrow 0.
\end{eqnarray*}

Let $v \in V$ and $w \in W$ be vectors  with $[T_xC] =[v + v^2 \otimes w]$. Then they define  vector subbundles $\mathcal V_0 \subset \mathcal V|_C$ and $\mathcal W_0 \subset \mathcal W|_C$. Set $\mathcal V_0^{(1)}:=\mathcal V|_C/\mathcal V_0$ and $\mathcal W_0^{(2)}:=\mathcal W|_C/\mathcal W_0$.
By Lemma \ref{lem:finer exact sequences}, we get the following exact sequences.
%
\begin{eqnarray*}
0 \rightarrow &\mathcal O(2)& \rightarrow\mathcal Q^{(0)} \rightarrow 0 \\
0 \rightarrow \mathcal V_0  \rightarrow &\mathcal O(1)^p& \rightarrow \mathcal Q^{(1)}  \rightarrow 0 \\
0 \rightarrow \mathcal V_0^{(1)}  \rightarrow &\mathcal O^r& \rightarrow \mathcal Q^{(2)}  \rightarrow 0\\
0 \rightarrow &\mathcal O(-1)^s& \rightarrow \mathcal Q^{(3)} \rightarrow 0
\end{eqnarray*}
with
\begin{eqnarray*}
  \mathcal Q^{(0)} &= &    \Sym^2 \mathcal V_0 \otimes \mathcal W_0  \\
\deg \mathcal Q^{(1)} &=& \deg \left(\mathcal V_0 \circ \mathcal V_0^{(1)}    \otimes \mathcal W_0\right)  \oplus \left( \Sym^2 \mathcal V_0 \otimes \mathcal W_0^{(1)} \right)  \\
\deg \mathcal Q^{(2)} &=& \deg \left(\Sym^2 \mathcal V_0^{(1)}  \otimes \mathcal W_0\right)  \oplus \left( \mathcal V_0 \circ  \mathcal V_0^{(1)}  \otimes \mathcal W _0^{(1)} \right)  \\
 \mathcal Q^{(3)} &= &     \Sym^2  \mathcal V _0^{(1)}  \otimes  \mathcal W_0^{(1)} .
\end{eqnarray*}
%
Furthermore,
\begin{enumerate}
\item[]
any direct summand of $\mathcal V_0^{(1)}$ has degree $\leq 0$, and
\item[]
any direct summand of $\Sym^2 \mathcal V_0^{(1)} \otimes \mathcal W_0^{(1)}$ has degree  $  -1$.
\end{enumerate}
The first statement follows from the third exact sequence, and the second statement follows from the last one.

\medskip
When $X$ is $(B_m, \alpha_{m-1}, \alpha_m)$, write $$\mathcal V_0  =\mathcal O(a_1), \mathcal V_0^{(1)}=\mathcal O(a_2), \mathcal W_0 =\mathcal O(b_1), \text{ and } \mathcal W_0^{(1)}=\mathcal O(b_2) \oplus \dots \oplus \mathcal O(b_{m-1}).$$
Then we have
\begin{eqnarray*}
2&=& 2a_1 + b_1 \\
m&=& a_1 +(a_1 +a_2 +b_1) + \sum_{i=2}^{m-1}(2a_1+b_i) =(2m-2)a_1 + a_2 + \sum_{i=1}^{m-1}b_i \\
0&=&a_2 + (2a_2 + b_1) + \sum_{i=2}^{m-1}(a_1 +a_2 + b_i) = (m-2)a_1 + (m+1) a_2 + \sum_{i=1}^{m-1} b_i \\
-m+2 &=& \sum_{i=2}^{m-1}(2a_2 + b_i) = 2(m-2) a_2+ \sum_{i=2}^{m-1} b_i.
\end{eqnarray*}
Therefore, $a_1=1$ and $a_2=0$ and $b_1=0$ and $\sum_{i=2}^{m-1} b_i = -m+2$.
Since the degree of a direct summand of $\Sym^2 \mathcal V_0^{(1)} \otimes \mathcal W_0^{(1)}$ is $  -1$, we have $b_i=-1$ for all $2 \leq i \leq m-1$.

Since the exact sequences
 \begin{eqnarray*}
 0 \rightarrow \mathcal V_0 \rightarrow &\mathcal V|_C& \rightarrow \mathcal V_0^{(1)} \rightarrow 0\\
 0 \rightarrow \mathcal W_0 \rightarrow &\mathcal W|_C& \rightarrow \mathcal W_0^{(1)} \rightarrow 0
 \end{eqnarray*}
split, we have the desired decompositions of
$\mathcal V|_C$ and $\mathcal W|_C$.


\medskip
When $X$ is $(F_4, \alpha_{2}, \alpha_3)$, write
$$\mathcal V_0  =\mathcal O(a_1), \mathcal V_0^{(1)}=   \mathcal O(a_2)\oplus \mathcal O(a_3), \mathcal W_0 =\mathcal O(b_1), \text{ and } \mathcal W_0^{(1)}= \mathcal O(b_2).$$
 Then we have
 \begin{eqnarray*}
 2&=& 2a_1 + b_1 \\
 4&=& a_1 + (a_1 +a_2 +b_1) +(a_1+a_3 +b_1) +(2a_1+b_2) \\
 &=& 5a_1 +(a_2 +a_3) +2 b_1 +b_2 \\
 0&=& a_2+a_3 + (3(a_2 +a_3) + 3 b_1)) + (a_1 +a_2 + b_2) + (a_1 +a_3+ b_2)\\
 &=& 2a_1 + 5(a_2 +a_3) +3 b_1 +2 b_2  \\
 -3&=& 3(a_2+a_3) +3b_2.
 \end{eqnarray*}
Therefore, $a_1=1$ and $b_1=0$ and $a_2+a_3=0$ and $b_2=-1$.
Since the degree of a direct summand of $\mathcal V_0^{(1)}$ is $\leq 0$, we get $a_2=a_3=0$.
By the same reason as in the case of $(B_m, \alpha_{m-1}, \alpha_m)$,  we have the desired decompositions of
$\mathcal V|_C$ and $\mathcal W|_C$.

This completes the proof of Lemma \ref{lem:splitting type of V}.
\end{proof}

\medskip
Now we continue the proof of Proposition \ref{extension to rationally connected open subset}.
Set
\begin{eqnarray*}
\mathcal E &:=&(\Sym^4\mathcal V)^{\perp} \otimes \wedge^2 \mathcal W = \Sym^2(\wedge^2 \mathcal V) \otimes \wedge^2 \mathcal W.
\end{eqnarray*}
By Lemma \ref{lem:splitting type of V}, we have
\begin{eqnarray*}
\mathcal E|_C
&=& \left\{\begin{array}{l}
 \mathcal O(2) \otimes (\mathcal O(-1)^{m-2} \oplus \mathcal O(-2)^t)   \text{ when } X \text{ is } (B_m, \alpha_{m-1}, \alpha_m) \\[5 pt]
\Sym^2 (\mathcal O(1)^2 \oplus \mathcal O) \otimes \mathcal O(-1)   \text{ when } X \text{ is } (F_4, \alpha_2, \alpha_3).
\end{array}\right.
\end{eqnarray*}
The Frobenius bracket $[\,,\,]:\wedge^2D|_C \rightarrow TM/D|_C$ induces a map
$$ \mathcal E |_C \rightarrow TM/D|_C, $$
which is an isomorphism onto its image for generic $y \in C$ by Proposition \ref{fundamental graded Lie algebra of vmrt} (1).

When $X$ is $(B_m, \alpha_{m-1}, \alpha_m)$, since $\mathcal E|_C$ and $TM/D|_C$ have the same splitting type, $\mathcal E|_C \rightarrow TM/D|_C$ is an isomorphism at any point of $C$.

When $X$ is $(F_4, \alpha_2, \alpha_3)$, $\mathcal E|_C \rightarrow TM/D|_C$ is an isomorphism onto $D^{-2}/D|_C$ at any point of $C$.
Now we have
$$0 \rightarrow D^{-2}/D \rightarrow TM/D \rightarrow TM/D^{-2} \rightarrow 0.$$
From $D^{-2}/D|_C=\mathcal O(1)^3 \oplus \mathcal O \oplus \mathcal O(-1)$ and $TM/D|_C=\mathcal O(1)^3 \oplus \mathcal O^5$, it follows that $TM/D^{-2}|_C= \mathcal O(1) \oplus \mathcal O$.

Set
\begin{eqnarray*}
 \mathcal E^{-3}&:=&\Sym^2(\wedge^3\mathcal V) \otimes (\mathcal W \otimes \wedge^2 \mathcal W)
\end{eqnarray*}
Then $\mathcal E^{-3}|_C = \mathcal O(1) \oplus \mathcal O$.
The Frobenius  bracket $[\,,\,]:D \wedge   (D^{-2}/D) \rightarrow  TM /D^{-2}$ induces a map

$$\mathcal E^{-3}|_C \rightarrow  TM  /D^{-2}|_{C} $$
which is an isomorphism for generic $y \in C$ by Proposition \ref{fundamental graded Lie algebra of vmrt} (1). Hence it  is an isomorphicm at any point of $C$ because  $TM/D^{-2}|_C =\mathcal O(1) \oplus \mathcal O$   has same splitting type as $\mathcal E^{-3}|_C$,

Consequently, $\Symb_y(D)$ is isomorphic to $\mathfrak m$ for any $y \in C$.


In sum, there exists a Zariski open subset $M'\subset M^0\subset M$ such that a general member  of $\mathcal K$ lies on $M^0$ and the varieties of minimal rational tangents $\mathcal C|_{M^0}\subset \mathbb P D|_{M^0}$ defines an ${\bf S}$-structure on $(M^0, D|_{M^0})$ of type $\frak m$.
%By Proposition \ref{S-str and G_0-str horospherical case}, % and Proposition \ref{fundamental graded Lie algebra of vmrt} (2),
Proposition \ref{S-str and G_0-str} together with Proposition \ref{fundamental graded Lie algebra of vmrt} implies that
there exist $G_0$-structure on $(M^0, D|_{M^0})$ corresponding to the ${\bf S}$-structure, where $G_0=G(\widehat{\bf S})$. This completes the proof of Proposition  \ref{extension to rationally connected open subset}.
\end{proof}


\begin{proof} [Proof of Theorem \ref{main results}] By Proposition \ref{extension to rationally connected open subset}, there is a Zariski open subset $M^0\subset M$ such that a general member  of $\mathcal K$ lies on $M^0$, satisfying  the following properties:
\begin{enumerate}
\item $D|_{M^0}$ is of type $\frak m$;
\item $\mathcal C|_{M^0}\subset \mathbb P D|_{M^0}$ defines an ${\bf S}$-structure on $(M^0, D|_{M^0})$ and there corresponds to a $G_0$-structure $\mathscr P$ on $(M^0, D|_{M^0})$, where $G_0=G(\widehat{\bf S})$.
\end{enumerate}
%
Define a vector bundle $\mathcal H^2_k$ on $M^0$  by  $\mathcal H^2_k:=\mathscr P \times _{G_0} H^2(\frak m, \frak g)_k$.

\begin{lemma} \label{lem:section vanishes} %We claim that
  $H^0(M^0, \mathcal H^2_k)$ is zero for all $k \geq 1$.
\end{lemma}


\begin{proof}
By Proposition \ref{prop:H2 vanishing}, it suffices to show that the followings;
\begin{enumerate}
\item $H^0(M^0, \wedge ^2 D^* \otimes \mathcal V)=0$ when $(\frak m, \frak g_0)$ is of type $(B_m, \alpha_{m-1}, \alpha_m )$ for $m>3$ or of type $(F_4, \alpha_2, \alpha_3 )$; 
\item $H^0(M^0, \wedge ^2 D^* \otimes D)=0$ when $(\frak m, \frak g_0)$ is of type $(B_m, \alpha_{m-1}, \alpha_m )$ for $m=3$;
\item  $H^0(M^0, (\wedge^2 D \otimes D)^* \otimes \mathcal V)=0$ when $(\frak m, \frak g_0)$ is of type $(B_m, \alpha_{m-1}, \alpha_m )$ for $m \geq 3$;
%\item $H^0(M^0, \wedge ^2 D^* \otimes \mathcal V)=0$ when $(\frak m, \frak g_0)$ is of type $(F_4, \alpha_2, \alpha_3 )$;
\end{enumerate}

We will adapt an argument similar to the proof of Proposition 6.2 of \cite{HH}. Assume that $(\frak m, \frak g_0)$ is either of type $(B_m, \alpha_{m-1}, \alpha_m )$ for $m > 3$ or of type $(F_4, \alpha_2, \alpha_3 )$. Let   $\varphi :\wedge^2 D \rightarrow \mathcal V$  be a nontrivial  vector bundle map.    For  $x \in M^0$ and $\beta \in D_x$ with $[\beta]\in \mathcal C_x$, take $C$ to be a member of $\mathcal K$ passing through $x$ with $[T_xC] =[\beta]$. Then $TC \wedge D|_C $  is decomposed as a sum of $\mathcal O(a)$'s with $a \geq 1$. Thus $\varphi|_C$ maps $TC \wedge D|_C $ into the $\mathcal O(1)$-factor of $\mathcal V|_C$, the intersection of $\mathcal V_x $ with $T_{\beta}\widehat{\mathcal C}_x$. Applying this argument to a general $[\beta] \in \mathcal C_x$, we see that the image of $\varphi_x$ is contained in the intersection
\begin{equation*}
    \cap_{[\beta] \in \mathcal C_x}(T_{ \beta }\widehat{\mathcal C}_x \cap \mathcal V_x),
\end{equation*}
which is $G_0$-invariant and degenerate in $\mathcal V_x$, contradicting to the fact that $\mathcal V$ is an irreducible $G_0$-bundle. Therefore, $\varphi_x$ is zero.

When $(\frak m, \frak g_0)$ is of type $(B_3, \alpha_{2}, \alpha_3 )$, let   $\varphi :\wedge^2 D \rightarrow D$  be a nontrivial  vector bundle map. As the same argument as above, $\varphi|_C$ maps $TC \wedge D|_C $ into the $\mathcal O(a)$-factors of $D|_C$ with $a \geq 1$, the $T_{\beta}\widehat{\mathcal C}_x$ for a general $[\beta] \in \mathcal C_x$. Hence, the image of $\varphi_x$ is contained in $\cap_{[\beta] \in \mathcal C_x} T_{ v }\widehat{\mathcal C}_x$. We refer Lemma \ref{lem:finer exact sequences} and let $\{ v_1, v_2 \}$ be a basis of $V$ and $w \in W$. If we consider $[\beta_1]$ and $[\beta_2]$ in ${\mathcal C}_x$ which are correspond to the vectors $v_1+v^2_1 \otimes w$ and $v_2+v^2_2\otimes w$ in $\bf U$, then $\cap_{[\beta] \in \mathcal C_x} T_{ v }\widehat{\mathcal C}_x$ should be contained in the complement subspace of $\mathcal V_x$ in $D_x$ corresponding to $\Sym^2 V \otimes W$. If the image of $\varphi_x$ is nonzero, then it should contain a vector in $\mathcal V_x$ by $U_0$ part of $G_0$-action, which is a contradiction. Therefore, $\varphi_x$ is zero.

Assume that $(\frak m, \frak g_0)$ is of type $(B_m, \alpha_{m-1}, \alpha_m )$ for $m \geq 3$. Let $\psi: \wedge ^2 D \otimes D \rightarrow \mathcal V$ be a vector bundle map.
Then, by the same argument as above, for  $x \in M^0$ and $[\beta] \in \mathcal C_x$ with $[\beta]=[T_xC]$, $\psi_x$ maps $\beta \wedge T^{(2)}_{\beta}\widehat{\mathcal C}_x \otimes D_x$ into the intersection $\mathcal V_x \cap T_{\beta}\widehat{\mathcal C}_x$.
 Thus the restriction of $\psi_x$ to $(\span \{\beta \wedge T_{\beta}^{(2)}\mathcal C_x : \beta \in \mathcal C_x\})\otimes D_x$ is zero.
By Proposition \ref{kernel of frobenius braket}, $\span \{\beta \wedge T_{\beta}^{(2)}\mathcal C_x: \beta \in \mathcal C_x\}$ is the kernel of the Frobenius bracket $[\,,\,]:\wedge ^2D \rightarrow TM/D$ at each point $x\in M^0$, and thus $\psi $ induces a bundle map
$$\mathcal E \otimes D \rightarrow \mathcal V.$$

On the other hand,
for  $x \in M^0$ and $[\beta] \in \mathcal C_x$ with $[\beta] = [T_xC]$, we have $\mathcal E|_C =\mathcal O(1)^{m-2} \oplus \mathcal O^t$. Therefore, by the induced map $\mathcal E \otimes D \rightarrow \mathcal V$, the image from $\mathcal E|_C \otimes TC$ to $\mathcal V|_C=\mathcal O(1) \oplus \mathcal O$ is zero. Since $[\beta] \in \mathcal C_x$ span $D_x$, the map $\psi_x$ is zero.

This completes the proof of Lemma \ref{lem:section vanishes}.
 \end{proof}

By Proposition \ref{prop:prolongation methods} and  Lemma \ref{lem:section vanishes}, $G_0$-structure $\mathscr P$  on $(M^0, D|_{M^0})$ is locally equivalent to the standard one. By Proposition \ref{S-str and G_0-str} together with  Proposition \ref{fundamental graded Lie algebra of vmrt}, the ${\bf S}$-structure on $M^0$ defined by $\mathcal C(M)|_{M^0}$ is locally equivalent to the standard one. By Theorem \ref{Cartan-Fubini},  a local map preserving the varieties of minimal rational tangents can be extended to a global biholomorphism. Hence, $M$ is biholomorphic to $X$.
This completes the proof of Theorem \ref{main results}.
\end{proof}


\begin{thebibliography}{}

\bibitem{HS73} P. Hilton and U. Stammbach, \emph{On the differentials in the Lyndon-Hochschild-Serre spectral sequence}, Bulletin of the AMS, {\bf 79}, no. 4 (1973) 796--799


\bibitem{HS53} G. Hochschild and J.-P. Serre, \emph{Cohomology of group extensions}, Trans. Amer. Math.
Soc. 74 (1953), 110--134.

\bibitem{HH}J. Hong and J.-M. Hwang, Characterization of the rational homogeneous space associciated to a long simple root by its variety of minimal rational tangents, Algebracis Geometry in East Asia-Hanoi 2005, Advanced Studies in Pure Mathematics 50 (2008) 217-236.

\bibitem{HM} J. Hong and T. Morimoto, Prolongations, invariants, and fundamental identities of geometric structures, arXiv:2203.05182

\bibitem{Hw15} J.-M.  Hwang,  Deformation of the space of lines on the 5-dimensional hyperquadric. Preprint.

\bibitem{HL19} J.-M. Hwang and Q. Li, Characterizing symplectic Grassmannians by varieties of minimal rational tangents, J. Differential Geom. 119(2), (2021) 309-381

\bibitem{HL21} J.-M. Hwang and Q. Li, Recognizing the $G_2$-horospherical manifold of Picard number 1 by varieties of minimal rational tangents, preprint.

\bibitem{HM98} J.-M. Hwang and N. Mok, Rigidity of irreducible Hermitian symmetric spaces of the compact type under K\"ahler deformation, Invent. Math. 131 (1998), 393-418
\bibitem{HM99}J.-M.Hwang, Varieties of minimal rational tangents on uniruled projective manifolds, Several Complex Variables, MSRI Publications Volume 37 (1999) 351-389.
\bibitem{HM01}J.-M. Hwang and N. Mok, Cartan-Fubini type extension of holomorphic maps for Fano manifolds of Picard number 1, Journal Math. Pures Aool. 80(2001) 563-575.
\bibitem{HM02}J.-M. Hwang and N. Mok, Deformation rigidity of the rational homogeneous space associated to a long simple root, Annales Scientifiques de l’École Normale Supérieure. Vol. 351 No. 2 (2002) 73-184.

\bibitem{Kol} J.  Koll\'ar,  Rational curves on algebraic varieties. Vol. 32. Springer-Verlag, 1996.

\bibitem{Ki0} S.-Y. Kim, Geometric structures modeled after smooth projective horospherical varieties of Picard number one, thesis, Seoul National University, https://s-space.snu.ac.kr/bitstream/10371/121311/1/000000133390.pdf

\bibitem{Ki}S.-Y. Kim, Geometric structures modeled on smooth horospherical varieties of Picard number one, Transform. Groups {\bf 22} (2017) no. 2 361--386


\bibitem{Ko61} B. Kostant, \emph{Lie algebra cohomology and the generalized Borel-Weil Theorem} Ann. of Math. {\bf 74} no. 2 (1961) 329--387.

\bibitem{Mk08} N. Mok, Recognizing certain rational homogeneous manifolds of Picard number one from their varieties of minimal rational tangents,  Third International Congree of CHinese Mathematicians. Part 1,2, 41-61, AMS/IP Stud. Adv. Math., 42, pt. 1,2, Amer. Math. Soc., Providence, RI, 2008

\bibitem{Mo93} T. Morimoto, \emph{Geometric structures on filtered manifolds}, Hokkaido Mathematical Journal {\bf 22} (1993) 263--347

\bibitem{Pa} B. Pasquier, On some smooth projective two-orbit varieties with Picard number 1, Mathematische Annalen 344 (2009) 963-987.

\bibitem{Ta} N. Tanaka, \emph{On the equivalence problems associated with simple graded Lie algebras}, Hokkaido Math. J. {\bf 8} (1979) 23--84
\end{thebibliography}


\end{document} 